% Môn: Vật lý
% Lớp: 10
% Chương: Chương III. Động lực học
% Bài: L10C3 Bài 14. Định luật 1 Newton
% Số câu: 162
% App đọc file này trực tiếp — sửa rồi Commit trên GitHub, bấm Đồng bộ GitHub .tex

% L10C3 Bài 14. Định luật 1 Newton

% ===== Câu 1 =====
% ID: AIQ_L10C3_FULL_0163
% Mức: NB
\dangbt{Nhận biết quán tính và nội dung định luật I Newton}
\begin{ex}
% ID: AIQ_L10C3_FULL_0163
% Mức: NB
Một vật chịu hợp lực $0\,\text{N}$. Theo định luật $I$ Newton, gia tốc của vật bằng bao nhiêu (kết quả theo đơn vị m/s^2).
\choice
{\True 0}
{1}
{0}
{0}
\loigiai{
Áp dụng hệ thức phù hợp của bài học, thay số được $0\,\text{m}$ /s^2.
}
\end{ex}

% ===== Câu 2 =====
% ID: AIQ_L10C3_FULL_0164
% Mức: NB
\begin{ex}
% ID: AIQ_L10C3_FULL_0164
% Mức: NB
Một vật chịu hợp lực $0\,\text{N}$. Theo định luật $I$ Newton, gia tốc của vật bằng bao nhiêu (kết quả theo đơn vị m/s^2).
\choice
{1}
{\True 0}
{0}
{0}
\loigiai{
Áp dụng hệ thức phù hợp của bài học, thay số được $0\,\text{m}$ /s^2.
}
\end{ex}

% ===== Câu 3 =====
% ID: AIQ_L10C3_FULL_0165
% Mức: NB
\begin{ex}
% ID: AIQ_L10C3_FULL_0165
% Mức: NB
Một vật chịu hợp lực $0\,\text{N}$. Theo định luật $I$ Newton, gia tốc của vật bằng bao nhiêu (kết quả theo đơn vị m/s^2).
\choice
{0}
{1}
{\True 0}
{0}
\loigiai{
Áp dụng hệ thức phù hợp của bài học, thay số được $0\,\text{m}$ /s^2.
}
\end{ex}

% ===== Câu 4 =====
% ID: AIQ_L10C3_FULL_0166
% Mức: NB
\begin{ex}
% ID: AIQ_L10C3_FULL_0166
% Mức: NB
Một vật chịu hợp lực $0\,\text{N}$. Theo định luật $I$ Newton, gia tốc của vật bằng bao nhiêu (kết quả theo đơn vị m/s^2).
\choice
{0}
{1}
{0}
{\True 0}
\loigiai{
Áp dụng hệ thức phù hợp của bài học, thay số được $0\,\text{m}$ /s^2.
}
\end{ex}

% ===== Câu 5 =====
% ID: AIQ_L10C3_FULL_0167
% Mức: TH
\begin{ex}
% ID: AIQ_L10C3_FULL_0167
% Mức: TH
Một vật chịu hợp lực $0\,\text{N}$. Theo định luật $I$ Newton, gia tốc của vật bằng bao nhiêu (kết quả theo đơn vị m/s^2).
\choice
{\True 0}
{1}
{0}
{0}
\loigiai{
Áp dụng hệ thức phù hợp của bài học, thay số được $0\,\text{m}$ /s^2.
}
\end{ex}

% ===== Câu 6 =====
% ID: AIQ_L10C3_FULL_0168
% Mức: TH
\begin{ex}
% ID: AIQ_L10C3_FULL_0168
% Mức: TH
Một vật chịu hợp lực $0\,\text{N}$. Theo định luật $I$ Newton, gia tốc của vật bằng bao nhiêu (kết quả theo đơn vị m/s^2).
\choice
{1}
{\True 0}
{0}
{0}
\loigiai{
Áp dụng hệ thức phù hợp của bài học, thay số được $0\,\text{m}$ /s^2.
}
\end{ex}

% ===== Câu 7 =====
% ID: AIQ_L10C3_FULL_0169
% Mức: TH
\begin{ex}
% ID: AIQ_L10C3_FULL_0169
% Mức: TH
Một vật chịu hợp lực $0\,\text{N}$. Theo định luật $I$ Newton, gia tốc của vật bằng bao nhiêu (kết quả theo đơn vị m/s^2).
\choice
{0}
{1}
{\True 0}
{0}
\loigiai{
Áp dụng hệ thức phù hợp của bài học, thay số được $0\,\text{m}$ /s^2.
}
\end{ex}

% ===== Câu 8 =====
% ID: AIQ_L10C3_FULL_0170
% Mức: TH
\begin{ex}
% ID: AIQ_L10C3_FULL_0170
% Mức: TH
Một vật chịu hợp lực $0\,\text{N}$. Theo định luật $I$ Newton, gia tốc của vật bằng bao nhiêu (kết quả theo đơn vị m/s^2).
\choice
{0}
{1}
{0}
{\True 0}
\loigiai{
Áp dụng hệ thức phù hợp của bài học, thay số được $0\,\text{m}$ /s^2.
}
\end{ex}

% ===== Câu 9 =====
% ID: AIQ_L10C3_FULL_0171
% Mức: VD
\begin{ex}
% ID: AIQ_L10C3_FULL_0171
% Mức: VD
Một vật chịu hợp lực $0\,\text{N}$. Theo định luật $I$ Newton, gia tốc của vật bằng bao nhiêu (kết quả theo đơn vị m/s^2).
\choice
{\True 0}
{1}
{0}
{0}
\loigiai{
Áp dụng hệ thức phù hợp của bài học, thay số được $0\,\text{m}$ /s^2.
}
\end{ex}

% ===== Câu 10 =====
% ID: AIQ_L10C3_FULL_0172
% Mức: VD
\begin{ex}
% ID: AIQ_L10C3_FULL_0172
% Mức: VD
Một vật chịu hợp lực $0\,\text{N}$. Theo định luật $I$ Newton, gia tốc của vật bằng bao nhiêu (kết quả theo đơn vị m/s^2).
\choice
{1}
{\True 0}
{0}
{0}
\loigiai{
Áp dụng hệ thức phù hợp của bài học, thay số được $0\,\text{m}$ /s^2.
}
\end{ex}

% ===== Câu 11 =====
% ID: AIQ_L10C3_FULL_0173
% Mức: TH
\begin{ex}
% ID: AIQ_L10C3_FULL_0173
% Mức: TH
Một vật chịu hợp lực $0\,\text{N}$. Theo định luật $I$ Newton, gia tốc của vật bằng bao nhiêu (kết quả theo đơn vị m/s^2).
\choiceTF
{\True Giá trị cần tìm là $0\,\text{m}$ /s^2.}
{Giá trị cần tìm là $1\,\text{m}$ /s^2.}
{\True Phải xác định đầy đủ các lực hoặc đại lượng liên quan trước khi tính.}
{Có thể bỏ qua điều kiện áp dụng của công thức.}
\loigiai{
\begin{itemchoice}
\item (Đúng) Giá trị cần tìm là $0\,\text{m}$ /s^2. (Vì): kết quả $0\,\text{m}$ /s^2. b) Sai. c) Đúng. d) Sai vì phải kiểm tra điều kiện áp dụng
\item (Sai) Giá trị cần tìm là $1\,\text{m}$ /s^2.
\item (Đúng) Phải xác định đầy đủ các lực hoặc đại lượng liên quan trước khi tính.
\item (Sai) Có thể bỏ qua điều kiện áp dụng của công thức.
\end{itemchoice}
}
\end{ex}

% ===== Câu 12 =====
% ID: AIQ_L10C3_FULL_0174
% Mức: TH
\begin{ex}
% ID: AIQ_L10C3_FULL_0174
% Mức: TH
Một vật chịu hợp lực $0\,\text{N}$. Theo định luật $I$ Newton, gia tốc của vật bằng bao nhiêu (kết quả theo đơn vị m/s^2).
\choiceTF
{\True Giá trị cần tìm là $0\,\text{m}$ /s^2.}
{Giá trị cần tìm là $1\,\text{m}$ /s^2.}
{\True Phải xác định đầy đủ các lực hoặc đại lượng liên quan trước khi tính.}
{Có thể bỏ qua điều kiện áp dụng của công thức.}
\loigiai{
\begin{itemchoice}
\item (Đúng) Giá trị cần tìm là $0\,\text{m}$ /s^2. (Vì): kết quả $0\,\text{m}$ /s^2. b) Sai. c) Đúng. d) Sai vì phải kiểm tra điều kiện áp dụng
\item (Sai) Giá trị cần tìm là $1\,\text{m}$ /s^2.
\item (Đúng) Phải xác định đầy đủ các lực hoặc đại lượng liên quan trước khi tính.
\item (Sai) Có thể bỏ qua điều kiện áp dụng của công thức.
\end{itemchoice}
}
\end{ex}

% ===== Câu 13 =====
% ID: AIQ_L10C3_FULL_0175
% Mức: VD
\begin{ex}
% ID: AIQ_L10C3_FULL_0175
% Mức: VD
Một vật chịu hợp lực $0\,\text{N}$. Theo định luật $I$ Newton, gia tốc của vật bằng bao nhiêu (kết quả theo đơn vị m/s^2).
\choiceTF
{\True Giá trị cần tìm là $0\,\text{m}$ /s^2.}
{Giá trị cần tìm là $1\,\text{m}$ /s^2.}
{\True Phải xác định đầy đủ các lực hoặc đại lượng liên quan trước khi tính.}
{Có thể bỏ qua điều kiện áp dụng của công thức.}
\loigiai{
\begin{itemchoice}
\item (Đúng) Giá trị cần tìm là $0\,\text{m}$ /s^2. (Vì): kết quả $0\,\text{m}$ /s^2. b) Sai. c) Đúng. d) Sai vì phải kiểm tra điều kiện áp dụng
\item (Sai) Giá trị cần tìm là $1\,\text{m}$ /s^2.
\item (Đúng) Phải xác định đầy đủ các lực hoặc đại lượng liên quan trước khi tính.
\item (Sai) Có thể bỏ qua điều kiện áp dụng của công thức.
\end{itemchoice}
}
\end{ex}

% ===== Câu 14 =====
% ID: AIQ_L10C3_FULL_0176
% Mức: VD
\begin{ex}
% ID: AIQ_L10C3_FULL_0176
% Mức: VD
Một vật chịu hợp lực $0\,\text{N}$. Theo định luật $I$ Newton, gia tốc của vật bằng bao nhiêu (kết quả theo đơn vị m/s^2).
\choiceTF
{\True Giá trị cần tìm là $0\,\text{m}$ /s^2.}
{Giá trị cần tìm là $1\,\text{m}$ /s^2.}
{\True Phải xác định đầy đủ các lực hoặc đại lượng liên quan trước khi tính.}
{Có thể bỏ qua điều kiện áp dụng của công thức.}
\loigiai{
\begin{itemchoice}
\item (Đúng) Giá trị cần tìm là $0\,\text{m}$ /s^2. (Vì): kết quả $0\,\text{m}$ /s^2. b) Sai. c) Đúng. d) Sai vì phải kiểm tra điều kiện áp dụng
\item (Sai) Giá trị cần tìm là $1\,\text{m}$ /s^2.
\item (Đúng) Phải xác định đầy đủ các lực hoặc đại lượng liên quan trước khi tính.
\item (Sai) Có thể bỏ qua điều kiện áp dụng của công thức.
\end{itemchoice}
}
\end{ex}

% ===== Câu 15 =====
% ID: AIQ_L10C3_FULL_0177
% Mức: VD
\begin{ex}
% ID: AIQ_L10C3_FULL_0177
% Mức: VD
Một vật chịu hợp lực $0\,\text{N}$. Theo định luật $I$ Newton, gia tốc của vật bằng bao nhiêu (kết quả theo đơn vị m/s^2).
\choiceTF
{\True Giá trị cần tìm là $0\,\text{m}$ /s^2.}
{Giá trị cần tìm là $1\,\text{m}$ /s^2.}
{\True Phải xác định đầy đủ các lực hoặc đại lượng liên quan trước khi tính.}
{Có thể bỏ qua điều kiện áp dụng của công thức.}
\loigiai{
\begin{itemchoice}
\item (Đúng) Giá trị cần tìm là $0\,\text{m}$ /s^2. (Vì): kết quả $0\,\text{m}$ /s^2. b) Sai. c) Đúng. d) Sai vì phải kiểm tra điều kiện áp dụng
\item (Sai) Giá trị cần tìm là $1\,\text{m}$ /s^2.
\item (Đúng) Phải xác định đầy đủ các lực hoặc đại lượng liên quan trước khi tính.
\item (Sai) Có thể bỏ qua điều kiện áp dụng của công thức.
\end{itemchoice}
}
\end{ex}

% ===== Câu 16 =====
% ID: AIQ_L10C3_FULL_0178
% Mức: TH
\begin{ex}
% ID: AIQ_L10C3_FULL_0178
% Mức: TH
Một vật chịu hợp lực $0\,\text{N}$. Theo định luật $I$ Newton, gia tốc của vật bằng bao nhiêu (kết quả theo đơn vị m/s^2).
\shortans{0}
\loigiai{
Thay số vào hệ thức của bài học thu được $0\,\text{m}$ /s^2.
}
\end{ex}

% ===== Câu 17 =====
% ID: AIQ_L10C3_FULL_0179
% Mức: TH
\begin{ex}
% ID: AIQ_L10C3_FULL_0179
% Mức: TH
Một vật chịu hợp lực $0\,\text{N}$. Theo định luật $I$ Newton, gia tốc của vật bằng bao nhiêu (kết quả theo đơn vị m/s^2).
\shortans{0}
\loigiai{
Thay số vào hệ thức của bài học thu được $0\,\text{m}$ /s^2.
}
\end{ex}

% ===== Câu 18 =====
% ID: AIQ_L10C3_FULL_0180
% Mức: TH
\begin{ex}
% ID: AIQ_L10C3_FULL_0180
% Mức: TH
Một vật chịu hợp lực $0\,\text{N}$. Theo định luật $I$ Newton, gia tốc của vật bằng bao nhiêu (kết quả theo đơn vị m/s^2).
\shortans{0}
\loigiai{
Thay số vào hệ thức của bài học thu được $0\,\text{m}$ /s^2.
}
\end{ex}

% ===== Câu 19 =====
% ID: AIQ_L10C3_FULL_0181
% Mức: VD
\begin{ex}
% ID: AIQ_L10C3_FULL_0181
% Mức: VD
Một vật chịu hợp lực $0\,\text{N}$. Theo định luật $I$ Newton, gia tốc của vật bằng bao nhiêu (kết quả theo đơn vị m/s^2).
\shortans{0}
\loigiai{
Thay số vào hệ thức của bài học thu được $0\,\text{m}$ /s^2.
}
\end{ex}

% ===== Câu 20 =====
% ID: AIQ_L10C3_FULL_0182
% Mức: VD
\begin{ex}
% ID: AIQ_L10C3_FULL_0182
% Mức: VD
Một vật chịu hợp lực $0\,\text{N}$. Theo định luật $I$ Newton, gia tốc của vật bằng bao nhiêu (kết quả theo đơn vị m/s^2).
\shortans{0}
\loigiai{
Thay số vào hệ thức của bài học thu được $0\,\text{m}$ /s^2.
}
\end{ex}

% ===== Câu 21 =====
% ID: AIQ_L10C3_FULL_0183
% Mức: VD
\begin{ex}
% ID: AIQ_L10C3_FULL_0183
% Mức: VD
Một vật chịu hợp lực $0\,\text{N}$. Theo định luật $I$ Newton, gia tốc của vật bằng bao nhiêu (kết quả theo đơn vị m/s^2).
\shortans{0}
\loigiai{
Thay số vào hệ thức của bài học thu được $0\,\text{m}$ /s^2.
}
\end{ex}

% ===== Câu 22 =====
% ID: AIQ_L10C3_FULL_0184
% Mức: VD
\begin{bt}
% ID: AIQ_L10C3_FULL_0184
% Mức: VD
Một vật chịu hợp lực $0\,\text{N}$. Theo định luật $I$ Newton, gia tốc của vật bằng bao nhiêu (kết quả theo đơn vị m/s^2).
\loigiai{
Xác định các đại lượng đã cho, chọn hệ thức phù hợp và thay số. Kết quả: $0\,\text{m}$ /s^2.
}
\end{bt}

% ===== Câu 23 =====
% ID: AIQ_L10C3_FULL_0185
% Mức: VD
\begin{bt}
% ID: AIQ_L10C3_FULL_0185
% Mức: VD
Một vật chịu hợp lực $0\,\text{N}$. Theo định luật $I$ Newton, gia tốc của vật bằng bao nhiêu (kết quả theo đơn vị m/s^2).
\loigiai{
Xác định các đại lượng đã cho, chọn hệ thức phù hợp và thay số. Kết quả: $0\,\text{m}$ /s^2.
}
\end{bt}

% ===== Câu 24 =====
% ID: AIQ_L10C3_FULL_0186
% Mức: VD
\begin{bt}
% ID: AIQ_L10C3_FULL_0186
% Mức: VD
Một vật chịu hợp lực $0\,\text{N}$. Theo định luật $I$ Newton, gia tốc của vật bằng bao nhiêu (kết quả theo đơn vị m/s^2).
\loigiai{
Xác định các đại lượng đã cho, chọn hệ thức phù hợp và thay số. Kết quả: $0\,\text{m}$ /s^2.
}
\end{bt}

% ===== Câu 25 =====
% ID: AIQ_L10C3_FULL_0187
% Mức: VDC
\begin{bt}
% ID: AIQ_L10C3_FULL_0187
% Mức: VDC
Một vật chịu hợp lực $0\,\text{N}$. Theo định luật $I$ Newton, gia tốc của vật bằng bao nhiêu (kết quả theo đơn vị m/s^2).
\loigiai{
Xác định các đại lượng đã cho, chọn hệ thức phù hợp và thay số. Kết quả: $0\,\text{m}$ /s^2.
}
\end{bt}

% ===== Câu 26 =====
% ID: AIQ_L10C3_FULL_0188
% Mức: VDC
\begin{bt}
% ID: AIQ_L10C3_FULL_0188
% Mức: VDC
Một vật chịu hợp lực $0\,\text{N}$. Theo định luật $I$ Newton, gia tốc của vật bằng bao nhiêu (kết quả theo đơn vị m/s^2).
\loigiai{
Xác định các đại lượng đã cho, chọn hệ thức phù hợp và thay số. Kết quả: $0\,\text{m}$ /s^2.
}
\end{bt}

% ===== Câu 27 =====
% ID: AIQ_L10C3_FULL_0189
% Mức: VDC
\begin{bt}
% ID: AIQ_L10C3_FULL_0189
% Mức: VDC
Một vật chịu hợp lực $0\,\text{N}$. Theo định luật $I$ Newton, gia tốc của vật bằng bao nhiêu (kết quả theo đơn vị m/s^2).
\loigiai{
Xác định các đại lượng đã cho, chọn hệ thức phù hợp và thay số. Kết quả: $0\,\text{m}$ /s^2.
}
\end{bt}

% ===== Câu 28 =====
% ID: AIQ_L10C3_FULL_0190
% Mức: NB
\dangbt{Phân tích trạng thái đứng yên hoặc chuyển động thẳng đều}
\begin{ex}
% ID: AIQ_L10C3_FULL_0190
% Mức: NB
Một vật chịu hợp lực $0\,\text{N}$. Theo định luật $I$ Newton, gia tốc của vật bằng bao nhiêu (kết quả theo đơn vị m/s^2).
\choice
{\True 0}
{1}
{0}
{0}
\loigiai{
Áp dụng hệ thức phù hợp của bài học, thay số được $0\,\text{m}$ /s^2.
}
\end{ex}

% ===== Câu 29 =====
% ID: AIQ_L10C3_FULL_0191
% Mức: NB
\begin{ex}
% ID: AIQ_L10C3_FULL_0191
% Mức: NB
Một vật chịu hợp lực $0\,\text{N}$. Theo định luật $I$ Newton, gia tốc của vật bằng bao nhiêu (kết quả theo đơn vị m/s^2).
\choice
{1}
{\True 0}
{0}
{0}
\loigiai{
Áp dụng hệ thức phù hợp của bài học, thay số được $0\,\text{m}$ /s^2.
}
\end{ex}

% ===== Câu 30 =====
% ID: AIQ_L10C3_FULL_0192
% Mức: NB
\begin{ex}
% ID: AIQ_L10C3_FULL_0192
% Mức: NB
Một vật chịu hợp lực $0\,\text{N}$. Theo định luật $I$ Newton, gia tốc của vật bằng bao nhiêu (kết quả theo đơn vị m/s^2).
\choice
{0}
{1}
{\True 0}
{0}
\loigiai{
Áp dụng hệ thức phù hợp của bài học, thay số được $0\,\text{m}$ /s^2.
}
\end{ex}

% ===== Câu 31 =====
% ID: AIQ_L10C3_FULL_0193
% Mức: NB
\begin{ex}
% ID: AIQ_L10C3_FULL_0193
% Mức: NB
Một vật chịu hợp lực $0\,\text{N}$. Theo định luật $I$ Newton, gia tốc của vật bằng bao nhiêu (kết quả theo đơn vị m/s^2).
\choice
{0}
{1}
{0}
{\True 0}
\loigiai{
Áp dụng hệ thức phù hợp của bài học, thay số được $0\,\text{m}$ /s^2.
}
\end{ex}

% ===== Câu 32 =====
% ID: AIQ_L10C3_FULL_0194
% Mức: TH
\begin{ex}
% ID: AIQ_L10C3_FULL_0194
% Mức: TH
Một vật chịu hợp lực $0\,\text{N}$. Theo định luật $I$ Newton, gia tốc của vật bằng bao nhiêu (kết quả theo đơn vị m/s^2).
\choice
{\True 0}
{1}
{0}
{0}
\loigiai{
Áp dụng hệ thức phù hợp của bài học, thay số được $0\,\text{m}$ /s^2.
}
\end{ex}

% ===== Câu 33 =====
% ID: AIQ_L10C3_FULL_0195
% Mức: TH
\begin{ex}
% ID: AIQ_L10C3_FULL_0195
% Mức: TH
Một vật chịu hợp lực $0\,\text{N}$. Theo định luật $I$ Newton, gia tốc của vật bằng bao nhiêu (kết quả theo đơn vị m/s^2).
\choice
{1}
{\True 0}
{0}
{0}
\loigiai{
Áp dụng hệ thức phù hợp của bài học, thay số được $0\,\text{m}$ /s^2.
}
\end{ex}

% ===== Câu 34 =====
% ID: AIQ_L10C3_FULL_0196
% Mức: TH
\begin{ex}
% ID: AIQ_L10C3_FULL_0196
% Mức: TH
Một vật chịu hợp lực $0\,\text{N}$. Theo định luật $I$ Newton, gia tốc của vật bằng bao nhiêu (kết quả theo đơn vị m/s^2).
\choice
{0}
{1}
{\True 0}
{0}
\loigiai{
Áp dụng hệ thức phù hợp của bài học, thay số được $0\,\text{m}$ /s^2.
}
\end{ex}

% ===== Câu 35 =====
% ID: AIQ_L10C3_FULL_0197
% Mức: TH
\begin{ex}
% ID: AIQ_L10C3_FULL_0197
% Mức: TH
Một vật chịu hợp lực $0\,\text{N}$. Theo định luật $I$ Newton, gia tốc của vật bằng bao nhiêu (kết quả theo đơn vị m/s^2).
\choice
{0}
{1}
{0}
{\True 0}
\loigiai{
Áp dụng hệ thức phù hợp của bài học, thay số được $0\,\text{m}$ /s^2.
}
\end{ex}

% ===== Câu 36 =====
% ID: AIQ_L10C3_FULL_0198
% Mức: VD
\begin{ex}
% ID: AIQ_L10C3_FULL_0198
% Mức: VD
Một vật chịu hợp lực $0\,\text{N}$. Theo định luật $I$ Newton, gia tốc của vật bằng bao nhiêu (kết quả theo đơn vị m/s^2).
\choice
{\True 0}
{1}
{0}
{0}
\loigiai{
Áp dụng hệ thức phù hợp của bài học, thay số được $0\,\text{m}$ /s^2.
}
\end{ex}

% ===== Câu 37 =====
% ID: AIQ_L10C3_FULL_0199
% Mức: VD
\begin{ex}
% ID: AIQ_L10C3_FULL_0199
% Mức: VD
Một vật chịu hợp lực $0\,\text{N}$. Theo định luật $I$ Newton, gia tốc của vật bằng bao nhiêu (kết quả theo đơn vị m/s^2).
\choice
{1}
{\True 0}
{0}
{0}
\loigiai{
Áp dụng hệ thức phù hợp của bài học, thay số được $0\,\text{m}$ /s^2.
}
\end{ex}

% ===== Câu 38 =====
% ID: AIQ_L10C3_FULL_0200
% Mức: TH
\begin{ex}
% ID: AIQ_L10C3_FULL_0200
% Mức: TH
Một vật chịu hợp lực $0\,\text{N}$. Theo định luật $I$ Newton, gia tốc của vật bằng bao nhiêu (kết quả theo đơn vị m/s^2).
\choiceTF
{\True Giá trị cần tìm là $0\,\text{m}$ /s^2.}
{Giá trị cần tìm là $1\,\text{m}$ /s^2.}
{\True Phải xác định đầy đủ các lực hoặc đại lượng liên quan trước khi tính.}
{Có thể bỏ qua điều kiện áp dụng của công thức.}
\loigiai{
\begin{itemchoice}
\item (Đúng) Giá trị cần tìm là $0\,\text{m}$ /s^2. (Vì): kết quả $0\,\text{m}$ /s^2. b) Sai. c) Đúng. d) Sai vì phải kiểm tra điều kiện áp dụng
\item (Sai) Giá trị cần tìm là $1\,\text{m}$ /s^2.
\item (Đúng) Phải xác định đầy đủ các lực hoặc đại lượng liên quan trước khi tính.
\item (Sai) Có thể bỏ qua điều kiện áp dụng của công thức.
\end{itemchoice}
}
\end{ex}

% ===== Câu 39 =====
% ID: AIQ_L10C3_FULL_0201
% Mức: TH
\begin{ex}
% ID: AIQ_L10C3_FULL_0201
% Mức: TH
Một vật chịu hợp lực $0\,\text{N}$. Theo định luật $I$ Newton, gia tốc của vật bằng bao nhiêu (kết quả theo đơn vị m/s^2).
\choiceTF
{\True Giá trị cần tìm là $0\,\text{m}$ /s^2.}
{Giá trị cần tìm là $1\,\text{m}$ /s^2.}
{\True Phải xác định đầy đủ các lực hoặc đại lượng liên quan trước khi tính.}
{Có thể bỏ qua điều kiện áp dụng của công thức.}
\loigiai{
\begin{itemchoice}
\item (Đúng) Giá trị cần tìm là $0\,\text{m}$ /s^2. (Vì): kết quả $0\,\text{m}$ /s^2. b) Sai. c) Đúng. d) Sai vì phải kiểm tra điều kiện áp dụng
\item (Sai) Giá trị cần tìm là $1\,\text{m}$ /s^2.
\item (Đúng) Phải xác định đầy đủ các lực hoặc đại lượng liên quan trước khi tính.
\item (Sai) Có thể bỏ qua điều kiện áp dụng của công thức.
\end{itemchoice}
}
\end{ex}

% ===== Câu 40 =====
% ID: AIQ_L10C3_FULL_0202
% Mức: VD
\begin{ex}
% ID: AIQ_L10C3_FULL_0202
% Mức: VD
Một vật chịu hợp lực $0\,\text{N}$. Theo định luật $I$ Newton, gia tốc của vật bằng bao nhiêu (kết quả theo đơn vị m/s^2).
\choiceTF
{\True Giá trị cần tìm là $0\,\text{m}$ /s^2.}
{Giá trị cần tìm là $1\,\text{m}$ /s^2.}
{\True Phải xác định đầy đủ các lực hoặc đại lượng liên quan trước khi tính.}
{Có thể bỏ qua điều kiện áp dụng của công thức.}
\loigiai{
\begin{itemchoice}
\item (Đúng) Giá trị cần tìm là $0\,\text{m}$ /s^2. (Vì): kết quả $0\,\text{m}$ /s^2. b) Sai. c) Đúng. d) Sai vì phải kiểm tra điều kiện áp dụng
\item (Sai) Giá trị cần tìm là $1\,\text{m}$ /s^2.
\item (Đúng) Phải xác định đầy đủ các lực hoặc đại lượng liên quan trước khi tính.
\item (Sai) Có thể bỏ qua điều kiện áp dụng của công thức.
\end{itemchoice}
}
\end{ex}

% ===== Câu 41 =====
% ID: AIQ_L10C3_FULL_0203
% Mức: VD
\begin{ex}
% ID: AIQ_L10C3_FULL_0203
% Mức: VD
Một vật chịu hợp lực $0\,\text{N}$. Theo định luật $I$ Newton, gia tốc của vật bằng bao nhiêu (kết quả theo đơn vị m/s^2).
\choiceTF
{\True Giá trị cần tìm là $0\,\text{m}$ /s^2.}
{Giá trị cần tìm là $1\,\text{m}$ /s^2.}
{\True Phải xác định đầy đủ các lực hoặc đại lượng liên quan trước khi tính.}
{Có thể bỏ qua điều kiện áp dụng của công thức.}
\loigiai{
\begin{itemchoice}
\item (Đúng) Giá trị cần tìm là $0\,\text{m}$ /s^2. (Vì): kết quả $0\,\text{m}$ /s^2. b) Sai. c) Đúng. d) Sai vì phải kiểm tra điều kiện áp dụng
\item (Sai) Giá trị cần tìm là $1\,\text{m}$ /s^2.
\item (Đúng) Phải xác định đầy đủ các lực hoặc đại lượng liên quan trước khi tính.
\item (Sai) Có thể bỏ qua điều kiện áp dụng của công thức.
\end{itemchoice}
}
\end{ex}

% ===== Câu 42 =====
% ID: AIQ_L10C3_FULL_0204
% Mức: VD
\begin{ex}
% ID: AIQ_L10C3_FULL_0204
% Mức: VD
Một vật chịu hợp lực $0\,\text{N}$. Theo định luật $I$ Newton, gia tốc của vật bằng bao nhiêu (kết quả theo đơn vị m/s^2).
\choiceTF
{\True Giá trị cần tìm là $0\,\text{m}$ /s^2.}
{Giá trị cần tìm là $1\,\text{m}$ /s^2.}
{\True Phải xác định đầy đủ các lực hoặc đại lượng liên quan trước khi tính.}
{Có thể bỏ qua điều kiện áp dụng của công thức.}
\loigiai{
\begin{itemchoice}
\item (Đúng) Giá trị cần tìm là $0\,\text{m}$ /s^2. (Vì): kết quả $0\,\text{m}$ /s^2. b) Sai. c) Đúng. d) Sai vì phải kiểm tra điều kiện áp dụng
\item (Sai) Giá trị cần tìm là $1\,\text{m}$ /s^2.
\item (Đúng) Phải xác định đầy đủ các lực hoặc đại lượng liên quan trước khi tính.
\item (Sai) Có thể bỏ qua điều kiện áp dụng của công thức.
\end{itemchoice}
}
\end{ex}

% ===== Câu 43 =====
% ID: AIQ_L10C3_FULL_0205
% Mức: TH
\begin{ex}
% ID: AIQ_L10C3_FULL_0205
% Mức: TH
Một vật chịu hợp lực $0\,\text{N}$. Theo định luật $I$ Newton, gia tốc của vật bằng bao nhiêu (kết quả theo đơn vị m/s^2).
\shortans{0}
\loigiai{
Thay số vào hệ thức của bài học thu được $0\,\text{m}$ /s^2.
}
\end{ex}

% ===== Câu 44 =====
% ID: AIQ_L10C3_FULL_0206
% Mức: TH
\begin{ex}
% ID: AIQ_L10C3_FULL_0206
% Mức: TH
Một vật chịu hợp lực $0\,\text{N}$. Theo định luật $I$ Newton, gia tốc của vật bằng bao nhiêu (kết quả theo đơn vị m/s^2).
\shortans{0}
\loigiai{
Thay số vào hệ thức của bài học thu được $0\,\text{m}$ /s^2.
}
\end{ex}

% ===== Câu 45 =====
% ID: AIQ_L10C3_FULL_0207
% Mức: TH
\begin{ex}
% ID: AIQ_L10C3_FULL_0207
% Mức: TH
Một vật chịu hợp lực $0\,\text{N}$. Theo định luật $I$ Newton, gia tốc của vật bằng bao nhiêu (kết quả theo đơn vị m/s^2).
\shortans{0}
\loigiai{
Thay số vào hệ thức của bài học thu được $0\,\text{m}$ /s^2.
}
\end{ex}

% ===== Câu 46 =====
% ID: AIQ_L10C3_FULL_0208
% Mức: VD
\begin{ex}
% ID: AIQ_L10C3_FULL_0208
% Mức: VD
Một vật chịu hợp lực $0\,\text{N}$. Theo định luật $I$ Newton, gia tốc của vật bằng bao nhiêu (kết quả theo đơn vị m/s^2).
\shortans{0}
\loigiai{
Thay số vào hệ thức của bài học thu được $0\,\text{m}$ /s^2.
}
\end{ex}

% ===== Câu 47 =====
% ID: AIQ_L10C3_FULL_0209
% Mức: VD
\begin{ex}
% ID: AIQ_L10C3_FULL_0209
% Mức: VD
Một vật chịu hợp lực $0\,\text{N}$. Theo định luật $I$ Newton, gia tốc của vật bằng bao nhiêu (kết quả theo đơn vị m/s^2).
\shortans{0}
\loigiai{
Thay số vào hệ thức của bài học thu được $0\,\text{m}$ /s^2.
}
\end{ex}

% ===== Câu 48 =====
% ID: AIQ_L10C3_FULL_0210
% Mức: VD
\begin{ex}
% ID: AIQ_L10C3_FULL_0210
% Mức: VD
Một vật chịu hợp lực $0\,\text{N}$. Theo định luật $I$ Newton, gia tốc của vật bằng bao nhiêu (kết quả theo đơn vị m/s^2).
\shortans{0}
\loigiai{
Thay số vào hệ thức của bài học thu được $0\,\text{m}$ /s^2.
}
\end{ex}

% ===== Câu 49 =====
% ID: AIQ_L10C3_FULL_0211
% Mức: VD
\begin{bt}
% ID: AIQ_L10C3_FULL_0211
% Mức: VD
Một vật chịu hợp lực $0\,\text{N}$. Theo định luật $I$ Newton, gia tốc của vật bằng bao nhiêu (kết quả theo đơn vị m/s^2).
\loigiai{
Xác định các đại lượng đã cho, chọn hệ thức phù hợp và thay số. Kết quả: $0\,\text{m}$ /s^2.
}
\end{bt}

% ===== Câu 50 =====
% ID: AIQ_L10C3_FULL_0212
% Mức: VD
\begin{bt}
% ID: AIQ_L10C3_FULL_0212
% Mức: VD
Một vật chịu hợp lực $0\,\text{N}$. Theo định luật $I$ Newton, gia tốc của vật bằng bao nhiêu (kết quả theo đơn vị m/s^2).
\loigiai{
Xác định các đại lượng đã cho, chọn hệ thức phù hợp và thay số. Kết quả: $0\,\text{m}$ /s^2.
}
\end{bt}

% ===== Câu 51 =====
% ID: AIQ_L10C3_FULL_0213
% Mức: VD
\begin{bt}
% ID: AIQ_L10C3_FULL_0213
% Mức: VD
Một vật chịu hợp lực $0\,\text{N}$. Theo định luật $I$ Newton, gia tốc của vật bằng bao nhiêu (kết quả theo đơn vị m/s^2).
\loigiai{
Xác định các đại lượng đã cho, chọn hệ thức phù hợp và thay số. Kết quả: $0\,\text{m}$ /s^2.
}
\end{bt}

% ===== Câu 52 =====
% ID: AIQ_L10C3_FULL_0214
% Mức: VDC
\begin{bt}
% ID: AIQ_L10C3_FULL_0214
% Mức: VDC
Một vật chịu hợp lực $0\,\text{N}$. Theo định luật $I$ Newton, gia tốc của vật bằng bao nhiêu (kết quả theo đơn vị m/s^2).
\loigiai{
Xác định các đại lượng đã cho, chọn hệ thức phù hợp và thay số. Kết quả: $0\,\text{m}$ /s^2.
}
\end{bt}

% ===== Câu 53 =====
% ID: AIQ_L10C3_FULL_0215
% Mức: VDC
\begin{bt}
% ID: AIQ_L10C3_FULL_0215
% Mức: VDC
Một vật chịu hợp lực $0\,\text{N}$. Theo định luật $I$ Newton, gia tốc của vật bằng bao nhiêu (kết quả theo đơn vị m/s^2).
\loigiai{
Xác định các đại lượng đã cho, chọn hệ thức phù hợp và thay số. Kết quả: $0\,\text{m}$ /s^2.
}
\end{bt}

% ===== Câu 54 =====
% ID: AIQ_L10C3_FULL_0216
% Mức: VDC
\begin{bt}
% ID: AIQ_L10C3_FULL_0216
% Mức: VDC
Một vật chịu hợp lực $0\,\text{N}$. Theo định luật $I$ Newton, gia tốc của vật bằng bao nhiêu (kết quả theo đơn vị m/s^2).
\loigiai{
Xác định các đại lượng đã cho, chọn hệ thức phù hợp và thay số. Kết quả: $0\,\text{m}$ /s^2.
}
\end{bt}

% ===== Câu 55 =====
% ID: AIQ_L10C3_FULL_0217
% Mức: NB
\dangbt{Giải thích hiện tượng quán tính trong thực tiễn}
\begin{ex}
% ID: AIQ_L10C3_FULL_0217
% Mức: NB
Một vật chịu hợp lực $0\,\text{N}$. Theo định luật $I$ Newton, gia tốc của vật bằng bao nhiêu (kết quả theo đơn vị m/s^2).
\choice
{\True 0}
{1}
{0}
{0}
\loigiai{
Áp dụng hệ thức phù hợp của bài học, thay số được $0\,\text{m}$ /s^2.
}
\end{ex}

% ===== Câu 56 =====
% ID: AIQ_L10C3_FULL_0218
% Mức: NB
\begin{ex}
% ID: AIQ_L10C3_FULL_0218
% Mức: NB
Một vật chịu hợp lực $0\,\text{N}$. Theo định luật $I$ Newton, gia tốc của vật bằng bao nhiêu (kết quả theo đơn vị m/s^2).
\choice
{1}
{\True 0}
{0}
{0}
\loigiai{
Áp dụng hệ thức phù hợp của bài học, thay số được $0\,\text{m}$ /s^2.
}
\end{ex}

% ===== Câu 57 =====
% ID: AIQ_L10C3_FULL_0219
% Mức: NB
\begin{ex}
% ID: AIQ_L10C3_FULL_0219
% Mức: NB
Một vật chịu hợp lực $0\,\text{N}$. Theo định luật $I$ Newton, gia tốc của vật bằng bao nhiêu (kết quả theo đơn vị m/s^2).
\choice
{0}
{1}
{\True 0}
{0}
\loigiai{
Áp dụng hệ thức phù hợp của bài học, thay số được $0\,\text{m}$ /s^2.
}
\end{ex}

% ===== Câu 58 =====
% ID: AIQ_L10C3_FULL_0220
% Mức: NB
\begin{ex}
% ID: AIQ_L10C3_FULL_0220
% Mức: NB
Một vật chịu hợp lực $0\,\text{N}$. Theo định luật $I$ Newton, gia tốc của vật bằng bao nhiêu (kết quả theo đơn vị m/s^2).
\choice
{0}
{1}
{0}
{\True 0}
\loigiai{
Áp dụng hệ thức phù hợp của bài học, thay số được $0\,\text{m}$ /s^2.
}
\end{ex}

% ===== Câu 59 =====
% ID: AIQ_L10C3_FULL_0221
% Mức: TH
\begin{ex}
% ID: AIQ_L10C3_FULL_0221
% Mức: TH
Một vật chịu hợp lực $0\,\text{N}$. Theo định luật $I$ Newton, gia tốc của vật bằng bao nhiêu (kết quả theo đơn vị m/s^2).
\choice
{\True 0}
{1}
{0}
{0}
\loigiai{
Áp dụng hệ thức phù hợp của bài học, thay số được $0\,\text{m}$ /s^2.
}
\end{ex}

% ===== Câu 60 =====
% ID: AIQ_L10C3_FULL_0222
% Mức: TH
\begin{ex}
% ID: AIQ_L10C3_FULL_0222
% Mức: TH
Một vật chịu hợp lực $0\,\text{N}$. Theo định luật $I$ Newton, gia tốc của vật bằng bao nhiêu (kết quả theo đơn vị m/s^2).
\choice
{1}
{\True 0}
{0}
{0}
\loigiai{
Áp dụng hệ thức phù hợp của bài học, thay số được $0\,\text{m}$ /s^2.
}
\end{ex}

% ===== Câu 61 =====
% ID: AIQ_L10C3_FULL_0223
% Mức: TH
\begin{ex}
% ID: AIQ_L10C3_FULL_0223
% Mức: TH
Một vật chịu hợp lực $0\,\text{N}$. Theo định luật $I$ Newton, gia tốc của vật bằng bao nhiêu (kết quả theo đơn vị m/s^2).
\choice
{0}
{1}
{\True 0}
{0}
\loigiai{
Áp dụng hệ thức phù hợp của bài học, thay số được $0\,\text{m}$ /s^2.
}
\end{ex}

% ===== Câu 62 =====
% ID: AIQ_L10C3_FULL_0224
% Mức: TH
\begin{ex}
% ID: AIQ_L10C3_FULL_0224
% Mức: TH
Một vật chịu hợp lực $0\,\text{N}$. Theo định luật $I$ Newton, gia tốc của vật bằng bao nhiêu (kết quả theo đơn vị m/s^2).
\choice
{0}
{1}
{0}
{\True 0}
\loigiai{
Áp dụng hệ thức phù hợp của bài học, thay số được $0\,\text{m}$ /s^2.
}
\end{ex}

% ===== Câu 63 =====
% ID: AIQ_L10C3_FULL_0225
% Mức: VD
\begin{ex}
% ID: AIQ_L10C3_FULL_0225
% Mức: VD
Một vật chịu hợp lực $0\,\text{N}$. Theo định luật $I$ Newton, gia tốc của vật bằng bao nhiêu (kết quả theo đơn vị m/s^2).
\choice
{\True 0}
{1}
{0}
{0}
\loigiai{
Áp dụng hệ thức phù hợp của bài học, thay số được $0\,\text{m}$ /s^2.
}
\end{ex}

% ===== Câu 64 =====
% ID: AIQ_L10C3_FULL_0226
% Mức: VD
\begin{ex}
% ID: AIQ_L10C3_FULL_0226
% Mức: VD
Một vật chịu hợp lực $0\,\text{N}$. Theo định luật $I$ Newton, gia tốc của vật bằng bao nhiêu (kết quả theo đơn vị m/s^2).
\choice
{1}
{\True 0}
{0}
{0}
\loigiai{
Áp dụng hệ thức phù hợp của bài học, thay số được $0\,\text{m}$ /s^2.
}
\end{ex}

% ===== Câu 65 =====
% ID: AIQ_L10C3_FULL_0227
% Mức: TH
\begin{ex}
% ID: AIQ_L10C3_FULL_0227
% Mức: TH
Một vật chịu hợp lực $0\,\text{N}$. Theo định luật $I$ Newton, gia tốc của vật bằng bao nhiêu (kết quả theo đơn vị m/s^2).
\choiceTF
{\True Giá trị cần tìm là $0\,\text{m}$ /s^2.}
{Giá trị cần tìm là $1\,\text{m}$ /s^2.}
{\True Phải xác định đầy đủ các lực hoặc đại lượng liên quan trước khi tính.}
{Có thể bỏ qua điều kiện áp dụng của công thức.}
\loigiai{
\begin{itemchoice}
\item (Đúng) Giá trị cần tìm là $0\,\text{m}$ /s^2. (Vì): kết quả $0\,\text{m}$ /s^2. b) Sai. c) Đúng. d) Sai vì phải kiểm tra điều kiện áp dụng
\item (Sai) Giá trị cần tìm là $1\,\text{m}$ /s^2.
\item (Đúng) Phải xác định đầy đủ các lực hoặc đại lượng liên quan trước khi tính.
\item (Sai) Có thể bỏ qua điều kiện áp dụng của công thức.
\end{itemchoice}
}
\end{ex}

% ===== Câu 66 =====
% ID: AIQ_L10C3_FULL_0228
% Mức: TH
\begin{ex}
% ID: AIQ_L10C3_FULL_0228
% Mức: TH
Một vật chịu hợp lực $0\,\text{N}$. Theo định luật $I$ Newton, gia tốc của vật bằng bao nhiêu (kết quả theo đơn vị m/s^2).
\choiceTF
{\True Giá trị cần tìm là $0\,\text{m}$ /s^2.}
{Giá trị cần tìm là $1\,\text{m}$ /s^2.}
{\True Phải xác định đầy đủ các lực hoặc đại lượng liên quan trước khi tính.}
{Có thể bỏ qua điều kiện áp dụng của công thức.}
\loigiai{
\begin{itemchoice}
\item (Đúng) Giá trị cần tìm là $0\,\text{m}$ /s^2. (Vì): kết quả $0\,\text{m}$ /s^2. b) Sai. c) Đúng. d) Sai vì phải kiểm tra điều kiện áp dụng
\item (Sai) Giá trị cần tìm là $1\,\text{m}$ /s^2.
\item (Đúng) Phải xác định đầy đủ các lực hoặc đại lượng liên quan trước khi tính.
\item (Sai) Có thể bỏ qua điều kiện áp dụng của công thức.
\end{itemchoice}
}
\end{ex}

% ===== Câu 67 =====
% ID: AIQ_L10C3_FULL_0229
% Mức: VD
\begin{ex}
% ID: AIQ_L10C3_FULL_0229
% Mức: VD
Một vật chịu hợp lực $0\,\text{N}$. Theo định luật $I$ Newton, gia tốc của vật bằng bao nhiêu (kết quả theo đơn vị m/s^2).
\choiceTF
{\True Giá trị cần tìm là $0\,\text{m}$ /s^2.}
{Giá trị cần tìm là $1\,\text{m}$ /s^2.}
{\True Phải xác định đầy đủ các lực hoặc đại lượng liên quan trước khi tính.}
{Có thể bỏ qua điều kiện áp dụng của công thức.}
\loigiai{
\begin{itemchoice}
\item (Đúng) Giá trị cần tìm là $0\,\text{m}$ /s^2. (Vì): kết quả $0\,\text{m}$ /s^2. b) Sai. c) Đúng. d) Sai vì phải kiểm tra điều kiện áp dụng
\item (Sai) Giá trị cần tìm là $1\,\text{m}$ /s^2.
\item (Đúng) Phải xác định đầy đủ các lực hoặc đại lượng liên quan trước khi tính.
\item (Sai) Có thể bỏ qua điều kiện áp dụng của công thức.
\end{itemchoice}
}
\end{ex}

% ===== Câu 68 =====
% ID: AIQ_L10C3_FULL_0230
% Mức: VD
\begin{ex}
% ID: AIQ_L10C3_FULL_0230
% Mức: VD
Một vật chịu hợp lực $0\,\text{N}$. Theo định luật $I$ Newton, gia tốc của vật bằng bao nhiêu (kết quả theo đơn vị m/s^2).
\choiceTF
{\True Giá trị cần tìm là $0\,\text{m}$ /s^2.}
{Giá trị cần tìm là $1\,\text{m}$ /s^2.}
{\True Phải xác định đầy đủ các lực hoặc đại lượng liên quan trước khi tính.}
{Có thể bỏ qua điều kiện áp dụng của công thức.}
\loigiai{
\begin{itemchoice}
\item (Đúng) Giá trị cần tìm là $0\,\text{m}$ /s^2. (Vì): kết quả $0\,\text{m}$ /s^2. b) Sai. c) Đúng. d) Sai vì phải kiểm tra điều kiện áp dụng
\item (Sai) Giá trị cần tìm là $1\,\text{m}$ /s^2.
\item (Đúng) Phải xác định đầy đủ các lực hoặc đại lượng liên quan trước khi tính.
\item (Sai) Có thể bỏ qua điều kiện áp dụng của công thức.
\end{itemchoice}
}
\end{ex}

% ===== Câu 69 =====
% ID: AIQ_L10C3_FULL_0231
% Mức: VD
\begin{ex}
% ID: AIQ_L10C3_FULL_0231
% Mức: VD
Một vật chịu hợp lực $0\,\text{N}$. Theo định luật $I$ Newton, gia tốc của vật bằng bao nhiêu (kết quả theo đơn vị m/s^2).
\choiceTF
{\True Giá trị cần tìm là $0\,\text{m}$ /s^2.}
{Giá trị cần tìm là $1\,\text{m}$ /s^2.}
{\True Phải xác định đầy đủ các lực hoặc đại lượng liên quan trước khi tính.}
{Có thể bỏ qua điều kiện áp dụng của công thức.}
\loigiai{
\begin{itemchoice}
\item (Đúng) Giá trị cần tìm là $0\,\text{m}$ /s^2. (Vì): kết quả $0\,\text{m}$ /s^2. b) Sai. c) Đúng. d) Sai vì phải kiểm tra điều kiện áp dụng
\item (Sai) Giá trị cần tìm là $1\,\text{m}$ /s^2.
\item (Đúng) Phải xác định đầy đủ các lực hoặc đại lượng liên quan trước khi tính.
\item (Sai) Có thể bỏ qua điều kiện áp dụng của công thức.
\end{itemchoice}
}
\end{ex}

% ===== Câu 70 =====
% ID: AIQ_L10C3_FULL_0232
% Mức: TH
\begin{ex}
% ID: AIQ_L10C3_FULL_0232
% Mức: TH
Một vật chịu hợp lực $0\,\text{N}$. Theo định luật $I$ Newton, gia tốc của vật bằng bao nhiêu (kết quả theo đơn vị m/s^2).
\shortans{0}
\loigiai{
Thay số vào hệ thức của bài học thu được $0\,\text{m}$ /s^2.
}
\end{ex}

% ===== Câu 71 =====
% ID: AIQ_L10C3_FULL_0233
% Mức: TH
\begin{ex}
% ID: AIQ_L10C3_FULL_0233
% Mức: TH
Một vật chịu hợp lực $0\,\text{N}$. Theo định luật $I$ Newton, gia tốc của vật bằng bao nhiêu (kết quả theo đơn vị m/s^2).
\shortans{0}
\loigiai{
Thay số vào hệ thức của bài học thu được $0\,\text{m}$ /s^2.
}
\end{ex}

% ===== Câu 72 =====
% ID: AIQ_L10C3_FULL_0234
% Mức: TH
\begin{ex}
% ID: AIQ_L10C3_FULL_0234
% Mức: TH
Một vật chịu hợp lực $0\,\text{N}$. Theo định luật $I$ Newton, gia tốc của vật bằng bao nhiêu (kết quả theo đơn vị m/s^2).
\shortans{0}
\loigiai{
Thay số vào hệ thức của bài học thu được $0\,\text{m}$ /s^2.
}
\end{ex}

% ===== Câu 73 =====
% ID: AIQ_L10C3_FULL_0235
% Mức: VD
\begin{ex}
% ID: AIQ_L10C3_FULL_0235
% Mức: VD
Một vật chịu hợp lực $0\,\text{N}$. Theo định luật $I$ Newton, gia tốc của vật bằng bao nhiêu (kết quả theo đơn vị m/s^2).
\shortans{0}
\loigiai{
Thay số vào hệ thức của bài học thu được $0\,\text{m}$ /s^2.
}
\end{ex}

% ===== Câu 74 =====
% ID: AIQ_L10C3_FULL_0236
% Mức: VD
\begin{ex}
% ID: AIQ_L10C3_FULL_0236
% Mức: VD
Một vật chịu hợp lực $0\,\text{N}$. Theo định luật $I$ Newton, gia tốc của vật bằng bao nhiêu (kết quả theo đơn vị m/s^2).
\shortans{0}
\loigiai{
Thay số vào hệ thức của bài học thu được $0\,\text{m}$ /s^2.
}
\end{ex}

% ===== Câu 75 =====
% ID: AIQ_L10C3_FULL_0237
% Mức: VD
\begin{ex}
% ID: AIQ_L10C3_FULL_0237
% Mức: VD
Một vật chịu hợp lực $0\,\text{N}$. Theo định luật $I$ Newton, gia tốc của vật bằng bao nhiêu (kết quả theo đơn vị m/s^2).
\shortans{0}
\loigiai{
Thay số vào hệ thức của bài học thu được $0\,\text{m}$ /s^2.
}
\end{ex}

% ===== Câu 76 =====
% ID: AIQ_L10C3_FULL_0238
% Mức: VD
\begin{bt}
% ID: AIQ_L10C3_FULL_0238
% Mức: VD
Một vật chịu hợp lực $0\,\text{N}$. Theo định luật $I$ Newton, gia tốc của vật bằng bao nhiêu (kết quả theo đơn vị m/s^2).
\loigiai{
Xác định các đại lượng đã cho, chọn hệ thức phù hợp và thay số. Kết quả: $0\,\text{m}$ /s^2.
}
\end{bt}

% ===== Câu 77 =====
% ID: AIQ_L10C3_FULL_0239
% Mức: VD
\begin{bt}
% ID: AIQ_L10C3_FULL_0239
% Mức: VD
Một vật chịu hợp lực $0\,\text{N}$. Theo định luật $I$ Newton, gia tốc của vật bằng bao nhiêu (kết quả theo đơn vị m/s^2).
\loigiai{
Xác định các đại lượng đã cho, chọn hệ thức phù hợp và thay số. Kết quả: $0\,\text{m}$ /s^2.
}
\end{bt}

% ===== Câu 78 =====
% ID: AIQ_L10C3_FULL_0240
% Mức: VD
\begin{bt}
% ID: AIQ_L10C3_FULL_0240
% Mức: VD
Một vật chịu hợp lực $0\,\text{N}$. Theo định luật $I$ Newton, gia tốc của vật bằng bao nhiêu (kết quả theo đơn vị m/s^2).
\loigiai{
Xác định các đại lượng đã cho, chọn hệ thức phù hợp và thay số. Kết quả: $0\,\text{m}$ /s^2.
}
\end{bt}

% ===== Câu 79 =====
% ID: AIQ_L10C3_FULL_0241
% Mức: VDC
\begin{bt}
% ID: AIQ_L10C3_FULL_0241
% Mức: VDC
Một vật chịu hợp lực $0\,\text{N}$. Theo định luật $I$ Newton, gia tốc của vật bằng bao nhiêu (kết quả theo đơn vị m/s^2).
\loigiai{
Xác định các đại lượng đã cho, chọn hệ thức phù hợp và thay số. Kết quả: $0\,\text{m}$ /s^2.
}
\end{bt}

% ===== Câu 80 =====
% ID: AIQ_L10C3_FULL_0242
% Mức: VDC
\begin{bt}
% ID: AIQ_L10C3_FULL_0242
% Mức: VDC
Một vật chịu hợp lực $0\,\text{N}$. Theo định luật $I$ Newton, gia tốc của vật bằng bao nhiêu (kết quả theo đơn vị m/s^2).
\loigiai{
Xác định các đại lượng đã cho, chọn hệ thức phù hợp và thay số. Kết quả: $0\,\text{m}$ /s^2.
}
\end{bt}

% ===== Câu 81 =====
% ID: AIQ_L10C3_FULL_0243
% Mức: VDC
\begin{bt}
% ID: AIQ_L10C3_FULL_0243
% Mức: VDC
Một vật chịu hợp lực $0\,\text{N}$. Theo định luật $I$ Newton, gia tốc của vật bằng bao nhiêu (kết quả theo đơn vị m/s^2).
\loigiai{
Xác định các đại lượng đã cho, chọn hệ thức phù hợp và thay số. Kết quả: $0\,\text{m}$ /s^2.
}
\end{bt}

% ===== Câu 82 =====
% ID: AIQ_L10C3_FULL_0244
% Mức: NB
\dangbt{Xác định lực cân bằng từ định luật I Newton}
\begin{ex}
% ID: AIQ_L10C3_FULL_0244
% Mức: NB
Một vật chịu hợp lực $0\,\text{N}$. Theo định luật $I$ Newton, gia tốc của vật bằng bao nhiêu (kết quả theo đơn vị m/s^2).
\choice
{\True 0}
{1}
{0}
{0}
\loigiai{
Áp dụng hệ thức phù hợp của bài học, thay số được $0\,\text{m}$ /s^2.
}
\end{ex}

% ===== Câu 83 =====
% ID: AIQ_L10C3_FULL_0245
% Mức: NB
\begin{ex}
% ID: AIQ_L10C3_FULL_0245
% Mức: NB
Một vật chịu hợp lực $0\,\text{N}$. Theo định luật $I$ Newton, gia tốc của vật bằng bao nhiêu (kết quả theo đơn vị m/s^2).
\choice
{1}
{\True 0}
{0}
{0}
\loigiai{
Áp dụng hệ thức phù hợp của bài học, thay số được $0\,\text{m}$ /s^2.
}
\end{ex}

% ===== Câu 84 =====
% ID: AIQ_L10C3_FULL_0246
% Mức: NB
\begin{ex}
% ID: AIQ_L10C3_FULL_0246
% Mức: NB
Một vật chịu hợp lực $0\,\text{N}$. Theo định luật $I$ Newton, gia tốc của vật bằng bao nhiêu (kết quả theo đơn vị m/s^2).
\choice
{0}
{1}
{\True 0}
{0}
\loigiai{
Áp dụng hệ thức phù hợp của bài học, thay số được $0\,\text{m}$ /s^2.
}
\end{ex}

% ===== Câu 85 =====
% ID: AIQ_L10C3_FULL_0247
% Mức: NB
\begin{ex}
% ID: AIQ_L10C3_FULL_0247
% Mức: NB
Một vật chịu hợp lực $0\,\text{N}$. Theo định luật $I$ Newton, gia tốc của vật bằng bao nhiêu (kết quả theo đơn vị m/s^2).
\choice
{0}
{1}
{0}
{\True 0}
\loigiai{
Áp dụng hệ thức phù hợp của bài học, thay số được $0\,\text{m}$ /s^2.
}
\end{ex}

% ===== Câu 86 =====
% ID: AIQ_L10C3_FULL_0248
% Mức: TH
\begin{ex}
% ID: AIQ_L10C3_FULL_0248
% Mức: TH
Một vật chịu hợp lực $0\,\text{N}$. Theo định luật $I$ Newton, gia tốc của vật bằng bao nhiêu (kết quả theo đơn vị m/s^2).
\choice
{\True 0}
{1}
{0}
{0}
\loigiai{
Áp dụng hệ thức phù hợp của bài học, thay số được $0\,\text{m}$ /s^2.
}
\end{ex}

% ===== Câu 87 =====
% ID: AIQ_L10C3_FULL_0249
% Mức: TH
\begin{ex}
% ID: AIQ_L10C3_FULL_0249
% Mức: TH
Một vật chịu hợp lực $0\,\text{N}$. Theo định luật $I$ Newton, gia tốc của vật bằng bao nhiêu (kết quả theo đơn vị m/s^2).
\choice
{1}
{\True 0}
{0}
{0}
\loigiai{
Áp dụng hệ thức phù hợp của bài học, thay số được $0\,\text{m}$ /s^2.
}
\end{ex}

% ===== Câu 88 =====
% ID: AIQ_L10C3_FULL_0250
% Mức: TH
\begin{ex}
% ID: AIQ_L10C3_FULL_0250
% Mức: TH
Một vật chịu hợp lực $0\,\text{N}$. Theo định luật $I$ Newton, gia tốc của vật bằng bao nhiêu (kết quả theo đơn vị m/s^2).
\choice
{0}
{1}
{\True 0}
{0}
\loigiai{
Áp dụng hệ thức phù hợp của bài học, thay số được $0\,\text{m}$ /s^2.
}
\end{ex}

% ===== Câu 89 =====
% ID: AIQ_L10C3_FULL_0251
% Mức: TH
\begin{ex}
% ID: AIQ_L10C3_FULL_0251
% Mức: TH
Một vật chịu hợp lực $0\,\text{N}$. Theo định luật $I$ Newton, gia tốc của vật bằng bao nhiêu (kết quả theo đơn vị m/s^2).
\choice
{0}
{1}
{0}
{\True 0}
\loigiai{
Áp dụng hệ thức phù hợp của bài học, thay số được $0\,\text{m}$ /s^2.
}
\end{ex}

% ===== Câu 90 =====
% ID: AIQ_L10C3_FULL_0252
% Mức: VD
\begin{ex}
% ID: AIQ_L10C3_FULL_0252
% Mức: VD
Một vật chịu hợp lực $0\,\text{N}$. Theo định luật $I$ Newton, gia tốc của vật bằng bao nhiêu (kết quả theo đơn vị m/s^2).
\choice
{\True 0}
{1}
{0}
{0}
\loigiai{
Áp dụng hệ thức phù hợp của bài học, thay số được $0\,\text{m}$ /s^2.
}
\end{ex}

% ===== Câu 91 =====
% ID: AIQ_L10C3_FULL_0253
% Mức: VD
\begin{ex}
% ID: AIQ_L10C3_FULL_0253
% Mức: VD
Một vật chịu hợp lực $0\,\text{N}$. Theo định luật $I$ Newton, gia tốc của vật bằng bao nhiêu (kết quả theo đơn vị m/s^2).
\choice
{1}
{\True 0}
{0}
{0}
\loigiai{
Áp dụng hệ thức phù hợp của bài học, thay số được $0\,\text{m}$ /s^2.
}
\end{ex}

% ===== Câu 92 =====
% ID: AIQ_L10C3_FULL_0254
% Mức: TH
\begin{ex}
% ID: AIQ_L10C3_FULL_0254
% Mức: TH
Một vật chịu hợp lực $0\,\text{N}$. Theo định luật $I$ Newton, gia tốc của vật bằng bao nhiêu (kết quả theo đơn vị m/s^2).
\choiceTF
{\True Giá trị cần tìm là $0\,\text{m}$ /s^2.}
{Giá trị cần tìm là $1\,\text{m}$ /s^2.}
{\True Phải xác định đầy đủ các lực hoặc đại lượng liên quan trước khi tính.}
{Có thể bỏ qua điều kiện áp dụng của công thức.}
\loigiai{
\begin{itemchoice}
\item (Đúng) Giá trị cần tìm là $0\,\text{m}$ /s^2. (Vì): kết quả $0\,\text{m}$ /s^2. b) Sai. c) Đúng. d) Sai vì phải kiểm tra điều kiện áp dụng
\item (Sai) Giá trị cần tìm là $1\,\text{m}$ /s^2.
\item (Đúng) Phải xác định đầy đủ các lực hoặc đại lượng liên quan trước khi tính.
\item (Sai) Có thể bỏ qua điều kiện áp dụng của công thức.
\end{itemchoice}
}
\end{ex}

% ===== Câu 93 =====
% ID: AIQ_L10C3_FULL_0255
% Mức: TH
\begin{ex}
% ID: AIQ_L10C3_FULL_0255
% Mức: TH
Một vật chịu hợp lực $0\,\text{N}$. Theo định luật $I$ Newton, gia tốc của vật bằng bao nhiêu (kết quả theo đơn vị m/s^2).
\choiceTF
{\True Giá trị cần tìm là $0\,\text{m}$ /s^2.}
{Giá trị cần tìm là $1\,\text{m}$ /s^2.}
{\True Phải xác định đầy đủ các lực hoặc đại lượng liên quan trước khi tính.}
{Có thể bỏ qua điều kiện áp dụng của công thức.}
\loigiai{
\begin{itemchoice}
\item (Đúng) Giá trị cần tìm là $0\,\text{m}$ /s^2. (Vì): kết quả $0\,\text{m}$ /s^2. b) Sai. c) Đúng. d) Sai vì phải kiểm tra điều kiện áp dụng
\item (Sai) Giá trị cần tìm là $1\,\text{m}$ /s^2.
\item (Đúng) Phải xác định đầy đủ các lực hoặc đại lượng liên quan trước khi tính.
\item (Sai) Có thể bỏ qua điều kiện áp dụng của công thức.
\end{itemchoice}
}
\end{ex}

% ===== Câu 94 =====
% ID: AIQ_L10C3_FULL_0256
% Mức: VD
\begin{ex}
% ID: AIQ_L10C3_FULL_0256
% Mức: VD
Một vật chịu hợp lực $0\,\text{N}$. Theo định luật $I$ Newton, gia tốc của vật bằng bao nhiêu (kết quả theo đơn vị m/s^2).
\choiceTF
{\True Giá trị cần tìm là $0\,\text{m}$ /s^2.}
{Giá trị cần tìm là $1\,\text{m}$ /s^2.}
{\True Phải xác định đầy đủ các lực hoặc đại lượng liên quan trước khi tính.}
{Có thể bỏ qua điều kiện áp dụng của công thức.}
\loigiai{
\begin{itemchoice}
\item (Đúng) Giá trị cần tìm là $0\,\text{m}$ /s^2. (Vì): kết quả $0\,\text{m}$ /s^2. b) Sai. c) Đúng. d) Sai vì phải kiểm tra điều kiện áp dụng
\item (Sai) Giá trị cần tìm là $1\,\text{m}$ /s^2.
\item (Đúng) Phải xác định đầy đủ các lực hoặc đại lượng liên quan trước khi tính.
\item (Sai) Có thể bỏ qua điều kiện áp dụng của công thức.
\end{itemchoice}
}
\end{ex}

% ===== Câu 95 =====
% ID: AIQ_L10C3_FULL_0257
% Mức: VD
\begin{ex}
% ID: AIQ_L10C3_FULL_0257
% Mức: VD
Một vật chịu hợp lực $0\,\text{N}$. Theo định luật $I$ Newton, gia tốc của vật bằng bao nhiêu (kết quả theo đơn vị m/s^2).
\choiceTF
{\True Giá trị cần tìm là $0\,\text{m}$ /s^2.}
{Giá trị cần tìm là $1\,\text{m}$ /s^2.}
{\True Phải xác định đầy đủ các lực hoặc đại lượng liên quan trước khi tính.}
{Có thể bỏ qua điều kiện áp dụng của công thức.}
\loigiai{
\begin{itemchoice}
\item (Đúng) Giá trị cần tìm là $0\,\text{m}$ /s^2. (Vì): kết quả $0\,\text{m}$ /s^2. b) Sai. c) Đúng. d) Sai vì phải kiểm tra điều kiện áp dụng
\item (Sai) Giá trị cần tìm là $1\,\text{m}$ /s^2.
\item (Đúng) Phải xác định đầy đủ các lực hoặc đại lượng liên quan trước khi tính.
\item (Sai) Có thể bỏ qua điều kiện áp dụng của công thức.
\end{itemchoice}
}
\end{ex}

% ===== Câu 96 =====
% ID: AIQ_L10C3_FULL_0258
% Mức: VD
\begin{ex}
% ID: AIQ_L10C3_FULL_0258
% Mức: VD
Một vật chịu hợp lực $0\,\text{N}$. Theo định luật $I$ Newton, gia tốc của vật bằng bao nhiêu (kết quả theo đơn vị m/s^2).
\choiceTF
{\True Giá trị cần tìm là $0\,\text{m}$ /s^2.}
{Giá trị cần tìm là $1\,\text{m}$ /s^2.}
{\True Phải xác định đầy đủ các lực hoặc đại lượng liên quan trước khi tính.}
{Có thể bỏ qua điều kiện áp dụng của công thức.}
\loigiai{
\begin{itemchoice}
\item (Đúng) Giá trị cần tìm là $0\,\text{m}$ /s^2. (Vì): kết quả $0\,\text{m}$ /s^2. b) Sai. c) Đúng. d) Sai vì phải kiểm tra điều kiện áp dụng
\item (Sai) Giá trị cần tìm là $1\,\text{m}$ /s^2.
\item (Đúng) Phải xác định đầy đủ các lực hoặc đại lượng liên quan trước khi tính.
\item (Sai) Có thể bỏ qua điều kiện áp dụng của công thức.
\end{itemchoice}
}
\end{ex}

% ===== Câu 97 =====
% ID: AIQ_L10C3_FULL_0259
% Mức: TH
\begin{ex}
% ID: AIQ_L10C3_FULL_0259
% Mức: TH
Một vật chịu hợp lực $0\,\text{N}$. Theo định luật $I$ Newton, gia tốc của vật bằng bao nhiêu (kết quả theo đơn vị m/s^2).
\shortans{0}
\loigiai{
Thay số vào hệ thức của bài học thu được $0\,\text{m}$ /s^2.
}
\end{ex}

% ===== Câu 98 =====
% ID: AIQ_L10C3_FULL_0260
% Mức: TH
\begin{ex}
% ID: AIQ_L10C3_FULL_0260
% Mức: TH
Một vật chịu hợp lực $0\,\text{N}$. Theo định luật $I$ Newton, gia tốc của vật bằng bao nhiêu (kết quả theo đơn vị m/s^2).
\shortans{0}
\loigiai{
Thay số vào hệ thức của bài học thu được $0\,\text{m}$ /s^2.
}
\end{ex}

% ===== Câu 99 =====
% ID: AIQ_L10C3_FULL_0261
% Mức: TH
\begin{ex}
% ID: AIQ_L10C3_FULL_0261
% Mức: TH
Một vật chịu hợp lực $0\,\text{N}$. Theo định luật $I$ Newton, gia tốc của vật bằng bao nhiêu (kết quả theo đơn vị m/s^2).
\shortans{0}
\loigiai{
Thay số vào hệ thức của bài học thu được $0\,\text{m}$ /s^2.
}
\end{ex}

% ===== Câu 100 =====
% ID: AIQ_L10C3_FULL_0262
% Mức: VD
\begin{ex}
% ID: AIQ_L10C3_FULL_0262
% Mức: VD
Một vật chịu hợp lực $0\,\text{N}$. Theo định luật $I$ Newton, gia tốc của vật bằng bao nhiêu (kết quả theo đơn vị m/s^2).
\shortans{0}
\loigiai{
Thay số vào hệ thức của bài học thu được $0\,\text{m}$ /s^2.
}
\end{ex}

% ===== Câu 101 =====
% ID: AIQ_L10C3_FULL_0263
% Mức: VD
\begin{ex}
% ID: AIQ_L10C3_FULL_0263
% Mức: VD
Một vật chịu hợp lực $0\,\text{N}$. Theo định luật $I$ Newton, gia tốc của vật bằng bao nhiêu (kết quả theo đơn vị m/s^2).
\shortans{0}
\loigiai{
Thay số vào hệ thức của bài học thu được $0\,\text{m}$ /s^2.
}
\end{ex}

% ===== Câu 102 =====
% ID: AIQ_L10C3_FULL_0264
% Mức: VD
\begin{ex}
% ID: AIQ_L10C3_FULL_0264
% Mức: VD
Một vật chịu hợp lực $0\,\text{N}$. Theo định luật $I$ Newton, gia tốc của vật bằng bao nhiêu (kết quả theo đơn vị m/s^2).
\shortans{0}
\loigiai{
Thay số vào hệ thức của bài học thu được $0\,\text{m}$ /s^2.
}
\end{ex}

% ===== Câu 103 =====
% ID: AIQ_L10C3_FULL_0265
% Mức: VD
\begin{bt}
% ID: AIQ_L10C3_FULL_0265
% Mức: VD
Một vật chịu hợp lực $0\,\text{N}$. Theo định luật $I$ Newton, gia tốc của vật bằng bao nhiêu (kết quả theo đơn vị m/s^2).
\loigiai{
Xác định các đại lượng đã cho, chọn hệ thức phù hợp và thay số. Kết quả: $0\,\text{m}$ /s^2.
}
\end{bt}

% ===== Câu 104 =====
% ID: AIQ_L10C3_FULL_0266
% Mức: VD
\begin{bt}
% ID: AIQ_L10C3_FULL_0266
% Mức: VD
Một vật chịu hợp lực $0\,\text{N}$. Theo định luật $I$ Newton, gia tốc của vật bằng bao nhiêu (kết quả theo đơn vị m/s^2).
\loigiai{
Xác định các đại lượng đã cho, chọn hệ thức phù hợp và thay số. Kết quả: $0\,\text{m}$ /s^2.
}
\end{bt}

% ===== Câu 105 =====
% ID: AIQ_L10C3_FULL_0267
% Mức: VD
\begin{bt}
% ID: AIQ_L10C3_FULL_0267
% Mức: VD
Một vật chịu hợp lực $0\,\text{N}$. Theo định luật $I$ Newton, gia tốc của vật bằng bao nhiêu (kết quả theo đơn vị m/s^2).
\loigiai{
Xác định các đại lượng đã cho, chọn hệ thức phù hợp và thay số. Kết quả: $0\,\text{m}$ /s^2.
}
\end{bt}

% ===== Câu 106 =====
% ID: AIQ_L10C3_FULL_0268
% Mức: VDC
\begin{bt}
% ID: AIQ_L10C3_FULL_0268
% Mức: VDC
Một vật chịu hợp lực $0\,\text{N}$. Theo định luật $I$ Newton, gia tốc của vật bằng bao nhiêu (kết quả theo đơn vị m/s^2).
\loigiai{
Xác định các đại lượng đã cho, chọn hệ thức phù hợp và thay số. Kết quả: $0\,\text{m}$ /s^2.
}
\end{bt}

% ===== Câu 107 =====
% ID: AIQ_L10C3_FULL_0269
% Mức: VDC
\begin{bt}
% ID: AIQ_L10C3_FULL_0269
% Mức: VDC
Một vật chịu hợp lực $0\,\text{N}$. Theo định luật $I$ Newton, gia tốc của vật bằng bao nhiêu (kết quả theo đơn vị m/s^2).
\loigiai{
Xác định các đại lượng đã cho, chọn hệ thức phù hợp và thay số. Kết quả: $0\,\text{m}$ /s^2.
}
\end{bt}

% ===== Câu 108 =====
% ID: AIQ_L10C3_FULL_0270
% Mức: VDC
\begin{bt}
% ID: AIQ_L10C3_FULL_0270
% Mức: VDC
Một vật chịu hợp lực $0\,\text{N}$. Theo định luật $I$ Newton, gia tốc của vật bằng bao nhiêu (kết quả theo đơn vị m/s^2).
\loigiai{
Xác định các đại lượng đã cho, chọn hệ thức phù hợp và thay số. Kết quả: $0\,\text{m}$ /s^2.
}
\end{bt}

% ===== Câu 109 =====
% ID: AIQ_L10C3_FULL_0271
% Mức: NB
\dangbt{Khai thác đồ thị chuyển động khi hợp lực bằng không}
\begin{ex}
% ID: AIQ_L10C3_FULL_0271
% Mức: NB
Một vật chịu hợp lực $0\,\text{N}$. Theo định luật $I$ Newton, gia tốc của vật bằng bao nhiêu (kết quả theo đơn vị m/s^2).
\choice
{\True 0}
{1}
{0}
{0}
\loigiai{
Áp dụng hệ thức phù hợp của bài học, thay số được $0\,\text{m}$ /s^2.
}
\end{ex}

% ===== Câu 110 =====
% ID: AIQ_L10C3_FULL_0272
% Mức: NB
\begin{ex}
% ID: AIQ_L10C3_FULL_0272
% Mức: NB
Một vật chịu hợp lực $0\,\text{N}$. Theo định luật $I$ Newton, gia tốc của vật bằng bao nhiêu (kết quả theo đơn vị m/s^2).
\choice
{1}
{\True 0}
{0}
{0}
\loigiai{
Áp dụng hệ thức phù hợp của bài học, thay số được $0\,\text{m}$ /s^2.
}
\end{ex}

% ===== Câu 111 =====
% ID: AIQ_L10C3_FULL_0273
% Mức: NB
\begin{ex}
% ID: AIQ_L10C3_FULL_0273
% Mức: NB
Một vật chịu hợp lực $0\,\text{N}$. Theo định luật $I$ Newton, gia tốc của vật bằng bao nhiêu (kết quả theo đơn vị m/s^2).
\choice
{0}
{1}
{\True 0}
{0}
\loigiai{
Áp dụng hệ thức phù hợp của bài học, thay số được $0\,\text{m}$ /s^2.
}
\end{ex}

% ===== Câu 112 =====
% ID: AIQ_L10C3_FULL_0274
% Mức: NB
\begin{ex}
% ID: AIQ_L10C3_FULL_0274
% Mức: NB
Một vật chịu hợp lực $0\,\text{N}$. Theo định luật $I$ Newton, gia tốc của vật bằng bao nhiêu (kết quả theo đơn vị m/s^2).
\choice
{0}
{1}
{0}
{\True 0}
\loigiai{
Áp dụng hệ thức phù hợp của bài học, thay số được $0\,\text{m}$ /s^2.
}
\end{ex}

% ===== Câu 113 =====
% ID: AIQ_L10C3_FULL_0275
% Mức: TH
\begin{ex}
% ID: AIQ_L10C3_FULL_0275
% Mức: TH
Một vật chịu hợp lực $0\,\text{N}$. Theo định luật $I$ Newton, gia tốc của vật bằng bao nhiêu (kết quả theo đơn vị m/s^2).
\choice
{\True 0}
{1}
{0}
{0}
\loigiai{
Áp dụng hệ thức phù hợp của bài học, thay số được $0\,\text{m}$ /s^2.
}
\end{ex}

% ===== Câu 114 =====
% ID: AIQ_L10C3_FULL_0276
% Mức: TH
\begin{ex}
% ID: AIQ_L10C3_FULL_0276
% Mức: TH
Một vật chịu hợp lực $0\,\text{N}$. Theo định luật $I$ Newton, gia tốc của vật bằng bao nhiêu (kết quả theo đơn vị m/s^2).
\choice
{1}
{\True 0}
{0}
{0}
\loigiai{
Áp dụng hệ thức phù hợp của bài học, thay số được $0\,\text{m}$ /s^2.
}
\end{ex}

% ===== Câu 115 =====
% ID: AIQ_L10C3_FULL_0277
% Mức: TH
\begin{ex}
% ID: AIQ_L10C3_FULL_0277
% Mức: TH
Một vật chịu hợp lực $0\,\text{N}$. Theo định luật $I$ Newton, gia tốc của vật bằng bao nhiêu (kết quả theo đơn vị m/s^2).
\choice
{0}
{1}
{\True 0}
{0}
\loigiai{
Áp dụng hệ thức phù hợp của bài học, thay số được $0\,\text{m}$ /s^2.
}
\end{ex}

% ===== Câu 116 =====
% ID: AIQ_L10C3_FULL_0278
% Mức: TH
\begin{ex}
% ID: AIQ_L10C3_FULL_0278
% Mức: TH
Một vật chịu hợp lực $0\,\text{N}$. Theo định luật $I$ Newton, gia tốc của vật bằng bao nhiêu (kết quả theo đơn vị m/s^2).
\choice
{0}
{1}
{0}
{\True 0}
\loigiai{
Áp dụng hệ thức phù hợp của bài học, thay số được $0\,\text{m}$ /s^2.
}
\end{ex}

% ===== Câu 117 =====
% ID: AIQ_L10C3_FULL_0279
% Mức: VD
\begin{ex}
% ID: AIQ_L10C3_FULL_0279
% Mức: VD
Một vật chịu hợp lực $0\,\text{N}$. Theo định luật $I$ Newton, gia tốc của vật bằng bao nhiêu (kết quả theo đơn vị m/s^2).
\choice
{\True 0}
{1}
{0}
{0}
\loigiai{
Áp dụng hệ thức phù hợp của bài học, thay số được $0\,\text{m}$ /s^2.
}
\end{ex}

% ===== Câu 118 =====
% ID: AIQ_L10C3_FULL_0280
% Mức: VD
\begin{ex}
% ID: AIQ_L10C3_FULL_0280
% Mức: VD
Một vật chịu hợp lực $0\,\text{N}$. Theo định luật $I$ Newton, gia tốc của vật bằng bao nhiêu (kết quả theo đơn vị m/s^2).
\choice
{1}
{\True 0}
{0}
{0}
\loigiai{
Áp dụng hệ thức phù hợp của bài học, thay số được $0\,\text{m}$ /s^2.
}
\end{ex}

% ===== Câu 119 =====
% ID: AIQ_L10C3_FULL_0281
% Mức: TH
\begin{ex}
% ID: AIQ_L10C3_FULL_0281
% Mức: TH
Một vật chịu hợp lực $0\,\text{N}$. Theo định luật $I$ Newton, gia tốc của vật bằng bao nhiêu (kết quả theo đơn vị m/s^2).
\choiceTF
{\True Giá trị cần tìm là $0\,\text{m}$ /s^2.}
{Giá trị cần tìm là $1\,\text{m}$ /s^2.}
{\True Phải xác định đầy đủ các lực hoặc đại lượng liên quan trước khi tính.}
{Có thể bỏ qua điều kiện áp dụng của công thức.}
\loigiai{
\begin{itemchoice}
\item (Đúng) Giá trị cần tìm là $0\,\text{m}$ /s^2. (Vì): kết quả $0\,\text{m}$ /s^2. b) Sai. c) Đúng. d) Sai vì phải kiểm tra điều kiện áp dụng
\item (Sai) Giá trị cần tìm là $1\,\text{m}$ /s^2.
\item (Đúng) Phải xác định đầy đủ các lực hoặc đại lượng liên quan trước khi tính.
\item (Sai) Có thể bỏ qua điều kiện áp dụng của công thức.
\end{itemchoice}
}
\end{ex}

% ===== Câu 120 =====
% ID: AIQ_L10C3_FULL_0282
% Mức: TH
\begin{ex}
% ID: AIQ_L10C3_FULL_0282
% Mức: TH
Một vật chịu hợp lực $0\,\text{N}$. Theo định luật $I$ Newton, gia tốc của vật bằng bao nhiêu (kết quả theo đơn vị m/s^2).
\choiceTF
{\True Giá trị cần tìm là $0\,\text{m}$ /s^2.}
{Giá trị cần tìm là $1\,\text{m}$ /s^2.}
{\True Phải xác định đầy đủ các lực hoặc đại lượng liên quan trước khi tính.}
{Có thể bỏ qua điều kiện áp dụng của công thức.}
\loigiai{
\begin{itemchoice}
\item (Đúng) Giá trị cần tìm là $0\,\text{m}$ /s^2. (Vì): kết quả $0\,\text{m}$ /s^2. b) Sai. c) Đúng. d) Sai vì phải kiểm tra điều kiện áp dụng
\item (Sai) Giá trị cần tìm là $1\,\text{m}$ /s^2.
\item (Đúng) Phải xác định đầy đủ các lực hoặc đại lượng liên quan trước khi tính.
\item (Sai) Có thể bỏ qua điều kiện áp dụng của công thức.
\end{itemchoice}
}
\end{ex}

% ===== Câu 121 =====
% ID: AIQ_L10C3_FULL_0283
% Mức: VD
\begin{ex}
% ID: AIQ_L10C3_FULL_0283
% Mức: VD
Một vật chịu hợp lực $0\,\text{N}$. Theo định luật $I$ Newton, gia tốc của vật bằng bao nhiêu (kết quả theo đơn vị m/s^2).
\choiceTF
{\True Giá trị cần tìm là $0\,\text{m}$ /s^2.}
{Giá trị cần tìm là $1\,\text{m}$ /s^2.}
{\True Phải xác định đầy đủ các lực hoặc đại lượng liên quan trước khi tính.}
{Có thể bỏ qua điều kiện áp dụng của công thức.}
\loigiai{
\begin{itemchoice}
\item (Đúng) Giá trị cần tìm là $0\,\text{m}$ /s^2. (Vì): kết quả $0\,\text{m}$ /s^2. b) Sai. c) Đúng. d) Sai vì phải kiểm tra điều kiện áp dụng
\item (Sai) Giá trị cần tìm là $1\,\text{m}$ /s^2.
\item (Đúng) Phải xác định đầy đủ các lực hoặc đại lượng liên quan trước khi tính.
\item (Sai) Có thể bỏ qua điều kiện áp dụng của công thức.
\end{itemchoice}
}
\end{ex}

% ===== Câu 122 =====
% ID: AIQ_L10C3_FULL_0284
% Mức: VD
\begin{ex}
% ID: AIQ_L10C3_FULL_0284
% Mức: VD
Một vật chịu hợp lực $0\,\text{N}$. Theo định luật $I$ Newton, gia tốc của vật bằng bao nhiêu (kết quả theo đơn vị m/s^2).
\choiceTF
{\True Giá trị cần tìm là $0\,\text{m}$ /s^2.}
{Giá trị cần tìm là $1\,\text{m}$ /s^2.}
{\True Phải xác định đầy đủ các lực hoặc đại lượng liên quan trước khi tính.}
{Có thể bỏ qua điều kiện áp dụng của công thức.}
\loigiai{
\begin{itemchoice}
\item (Đúng) Giá trị cần tìm là $0\,\text{m}$ /s^2. (Vì): kết quả $0\,\text{m}$ /s^2. b) Sai. c) Đúng. d) Sai vì phải kiểm tra điều kiện áp dụng
\item (Sai) Giá trị cần tìm là $1\,\text{m}$ /s^2.
\item (Đúng) Phải xác định đầy đủ các lực hoặc đại lượng liên quan trước khi tính.
\item (Sai) Có thể bỏ qua điều kiện áp dụng của công thức.
\end{itemchoice}
}
\end{ex}

% ===== Câu 123 =====
% ID: AIQ_L10C3_FULL_0285
% Mức: VD
\begin{ex}
% ID: AIQ_L10C3_FULL_0285
% Mức: VD
Một vật chịu hợp lực $0\,\text{N}$. Theo định luật $I$ Newton, gia tốc của vật bằng bao nhiêu (kết quả theo đơn vị m/s^2).
\choiceTF
{\True Giá trị cần tìm là $0\,\text{m}$ /s^2.}
{Giá trị cần tìm là $1\,\text{m}$ /s^2.}
{\True Phải xác định đầy đủ các lực hoặc đại lượng liên quan trước khi tính.}
{Có thể bỏ qua điều kiện áp dụng của công thức.}
\loigiai{
\begin{itemchoice}
\item (Đúng) Giá trị cần tìm là $0\,\text{m}$ /s^2. (Vì): kết quả $0\,\text{m}$ /s^2. b) Sai. c) Đúng. d) Sai vì phải kiểm tra điều kiện áp dụng
\item (Sai) Giá trị cần tìm là $1\,\text{m}$ /s^2.
\item (Đúng) Phải xác định đầy đủ các lực hoặc đại lượng liên quan trước khi tính.
\item (Sai) Có thể bỏ qua điều kiện áp dụng của công thức.
\end{itemchoice}
}
\end{ex}

% ===== Câu 124 =====
% ID: AIQ_L10C3_FULL_0286
% Mức: TH
\begin{ex}
% ID: AIQ_L10C3_FULL_0286
% Mức: TH
Một vật chịu hợp lực $0\,\text{N}$. Theo định luật $I$ Newton, gia tốc của vật bằng bao nhiêu (kết quả theo đơn vị m/s^2).
\shortans{0}
\loigiai{
Thay số vào hệ thức của bài học thu được $0\,\text{m}$ /s^2.
}
\end{ex}

% ===== Câu 125 =====
% ID: AIQ_L10C3_FULL_0287
% Mức: TH
\begin{ex}
% ID: AIQ_L10C3_FULL_0287
% Mức: TH
Một vật chịu hợp lực $0\,\text{N}$. Theo định luật $I$ Newton, gia tốc của vật bằng bao nhiêu (kết quả theo đơn vị m/s^2).
\shortans{0}
\loigiai{
Thay số vào hệ thức của bài học thu được $0\,\text{m}$ /s^2.
}
\end{ex}

% ===== Câu 126 =====
% ID: AIQ_L10C3_FULL_0288
% Mức: TH
\begin{ex}
% ID: AIQ_L10C3_FULL_0288
% Mức: TH
Một vật chịu hợp lực $0\,\text{N}$. Theo định luật $I$ Newton, gia tốc của vật bằng bao nhiêu (kết quả theo đơn vị m/s^2).
\shortans{0}
\loigiai{
Thay số vào hệ thức của bài học thu được $0\,\text{m}$ /s^2.
}
\end{ex}

% ===== Câu 127 =====
% ID: AIQ_L10C3_FULL_0289
% Mức: VD
\begin{ex}
% ID: AIQ_L10C3_FULL_0289
% Mức: VD
Một vật chịu hợp lực $0\,\text{N}$. Theo định luật $I$ Newton, gia tốc của vật bằng bao nhiêu (kết quả theo đơn vị m/s^2).
\shortans{0}
\loigiai{
Thay số vào hệ thức của bài học thu được $0\,\text{m}$ /s^2.
}
\end{ex}

% ===== Câu 128 =====
% ID: AIQ_L10C3_FULL_0290
% Mức: VD
\begin{ex}
% ID: AIQ_L10C3_FULL_0290
% Mức: VD
Một vật chịu hợp lực $0\,\text{N}$. Theo định luật $I$ Newton, gia tốc của vật bằng bao nhiêu (kết quả theo đơn vị m/s^2).
\shortans{0}
\loigiai{
Thay số vào hệ thức của bài học thu được $0\,\text{m}$ /s^2.
}
\end{ex}

% ===== Câu 129 =====
% ID: AIQ_L10C3_FULL_0291
% Mức: VD
\begin{ex}
% ID: AIQ_L10C3_FULL_0291
% Mức: VD
Một vật chịu hợp lực $0\,\text{N}$. Theo định luật $I$ Newton, gia tốc của vật bằng bao nhiêu (kết quả theo đơn vị m/s^2).
\shortans{0}
\loigiai{
Thay số vào hệ thức của bài học thu được $0\,\text{m}$ /s^2.
}
\end{ex}

% ===== Câu 130 =====
% ID: AIQ_L10C3_FULL_0292
% Mức: VD
\begin{bt}
% ID: AIQ_L10C3_FULL_0292
% Mức: VD
Một vật chịu hợp lực $0\,\text{N}$. Theo định luật $I$ Newton, gia tốc của vật bằng bao nhiêu (kết quả theo đơn vị m/s^2).
\loigiai{
Xác định các đại lượng đã cho, chọn hệ thức phù hợp và thay số. Kết quả: $0\,\text{m}$ /s^2.
}
\end{bt}

% ===== Câu 131 =====
% ID: AIQ_L10C3_FULL_0293
% Mức: VD
\begin{bt}
% ID: AIQ_L10C3_FULL_0293
% Mức: VD
Một vật chịu hợp lực $0\,\text{N}$. Theo định luật $I$ Newton, gia tốc của vật bằng bao nhiêu (kết quả theo đơn vị m/s^2).
\loigiai{
Xác định các đại lượng đã cho, chọn hệ thức phù hợp và thay số. Kết quả: $0\,\text{m}$ /s^2.
}
\end{bt}

% ===== Câu 132 =====
% ID: AIQ_L10C3_FULL_0294
% Mức: VD
\begin{bt}
% ID: AIQ_L10C3_FULL_0294
% Mức: VD
Một vật chịu hợp lực $0\,\text{N}$. Theo định luật $I$ Newton, gia tốc của vật bằng bao nhiêu (kết quả theo đơn vị m/s^2).
\loigiai{
Xác định các đại lượng đã cho, chọn hệ thức phù hợp và thay số. Kết quả: $0\,\text{m}$ /s^2.
}
\end{bt}

% ===== Câu 133 =====
% ID: AIQ_L10C3_FULL_0295
% Mức: VDC
\begin{bt}
% ID: AIQ_L10C3_FULL_0295
% Mức: VDC
Một vật chịu hợp lực $0\,\text{N}$. Theo định luật $I$ Newton, gia tốc của vật bằng bao nhiêu (kết quả theo đơn vị m/s^2).
\loigiai{
Xác định các đại lượng đã cho, chọn hệ thức phù hợp và thay số. Kết quả: $0\,\text{m}$ /s^2.
}
\end{bt}

% ===== Câu 134 =====
% ID: AIQ_L10C3_FULL_0296
% Mức: VDC
\begin{bt}
% ID: AIQ_L10C3_FULL_0296
% Mức: VDC
Một vật chịu hợp lực $0\,\text{N}$. Theo định luật $I$ Newton, gia tốc của vật bằng bao nhiêu (kết quả theo đơn vị m/s^2).
\loigiai{
Xác định các đại lượng đã cho, chọn hệ thức phù hợp và thay số. Kết quả: $0\,\text{m}$ /s^2.
}
\end{bt}

% ===== Câu 135 =====
% ID: AIQ_L10C3_FULL_0297
% Mức: VDC
\begin{bt}
% ID: AIQ_L10C3_FULL_0297
% Mức: VDC
Một vật chịu hợp lực $0\,\text{N}$. Theo định luật $I$ Newton, gia tốc của vật bằng bao nhiêu (kết quả theo đơn vị m/s^2).
\loigiai{
Xác định các đại lượng đã cho, chọn hệ thức phù hợp và thay số. Kết quả: $0\,\text{m}$ /s^2.
}
\end{bt}

% ===== Câu 136 =====
% ID: AIQ_L10C3_FULL_0298
% Mức: NB
\dangbt{Bài toán tổng hợp về quán tính và cân bằng lực}
\begin{ex}
% ID: AIQ_L10C3_FULL_0298
% Mức: NB
Một vật chịu hợp lực $0\,\text{N}$. Theo định luật $I$ Newton, gia tốc của vật bằng bao nhiêu (kết quả theo đơn vị m/s^2).
\choice
{\True 0}
{1}
{0}
{0}
\loigiai{
Áp dụng hệ thức phù hợp của bài học, thay số được $0\,\text{m}$ /s^2.
}
\end{ex}

% ===== Câu 137 =====
% ID: AIQ_L10C3_FULL_0299
% Mức: NB
\begin{ex}
% ID: AIQ_L10C3_FULL_0299
% Mức: NB
Một vật chịu hợp lực $0\,\text{N}$. Theo định luật $I$ Newton, gia tốc của vật bằng bao nhiêu (kết quả theo đơn vị m/s^2).
\choice
{1}
{\True 0}
{0}
{0}
\loigiai{
Áp dụng hệ thức phù hợp của bài học, thay số được $0\,\text{m}$ /s^2.
}
\end{ex}

% ===== Câu 138 =====
% ID: AIQ_L10C3_FULL_0300
% Mức: NB
\begin{ex}
% ID: AIQ_L10C3_FULL_0300
% Mức: NB
Một vật chịu hợp lực $0\,\text{N}$. Theo định luật $I$ Newton, gia tốc của vật bằng bao nhiêu (kết quả theo đơn vị m/s^2).
\choice
{0}
{1}
{\True 0}
{0}
\loigiai{
Áp dụng hệ thức phù hợp của bài học, thay số được $0\,\text{m}$ /s^2.
}
\end{ex}

% ===== Câu 139 =====
% ID: AIQ_L10C3_FULL_0301
% Mức: NB
\begin{ex}
% ID: AIQ_L10C3_FULL_0301
% Mức: NB
Một vật chịu hợp lực $0\,\text{N}$. Theo định luật $I$ Newton, gia tốc của vật bằng bao nhiêu (kết quả theo đơn vị m/s^2).
\choice
{0}
{1}
{0}
{\True 0}
\loigiai{
Áp dụng hệ thức phù hợp của bài học, thay số được $0\,\text{m}$ /s^2.
}
\end{ex}

% ===== Câu 140 =====
% ID: AIQ_L10C3_FULL_0302
% Mức: TH
\begin{ex}
% ID: AIQ_L10C3_FULL_0302
% Mức: TH
Một vật chịu hợp lực $0\,\text{N}$. Theo định luật $I$ Newton, gia tốc của vật bằng bao nhiêu (kết quả theo đơn vị m/s^2).
\choice
{\True 0}
{1}
{0}
{0}
\loigiai{
Áp dụng hệ thức phù hợp của bài học, thay số được $0\,\text{m}$ /s^2.
}
\end{ex}

% ===== Câu 141 =====
% ID: AIQ_L10C3_FULL_0303
% Mức: TH
\begin{ex}
% ID: AIQ_L10C3_FULL_0303
% Mức: TH
Một vật chịu hợp lực $0\,\text{N}$. Theo định luật $I$ Newton, gia tốc của vật bằng bao nhiêu (kết quả theo đơn vị m/s^2).
\choice
{1}
{\True 0}
{0}
{0}
\loigiai{
Áp dụng hệ thức phù hợp của bài học, thay số được $0\,\text{m}$ /s^2.
}
\end{ex}

% ===== Câu 142 =====
% ID: AIQ_L10C3_FULL_0304
% Mức: TH
\begin{ex}
% ID: AIQ_L10C3_FULL_0304
% Mức: TH
Một vật chịu hợp lực $0\,\text{N}$. Theo định luật $I$ Newton, gia tốc của vật bằng bao nhiêu (kết quả theo đơn vị m/s^2).
\choice
{0}
{1}
{\True 0}
{0}
\loigiai{
Áp dụng hệ thức phù hợp của bài học, thay số được $0\,\text{m}$ /s^2.
}
\end{ex}

% ===== Câu 143 =====
% ID: AIQ_L10C3_FULL_0305
% Mức: TH
\begin{ex}
% ID: AIQ_L10C3_FULL_0305
% Mức: TH
Một vật chịu hợp lực $0\,\text{N}$. Theo định luật $I$ Newton, gia tốc của vật bằng bao nhiêu (kết quả theo đơn vị m/s^2).
\choice
{0}
{1}
{0}
{\True 0}
\loigiai{
Áp dụng hệ thức phù hợp của bài học, thay số được $0\,\text{m}$ /s^2.
}
\end{ex}

% ===== Câu 144 =====
% ID: AIQ_L10C3_FULL_0306
% Mức: VD
\begin{ex}
% ID: AIQ_L10C3_FULL_0306
% Mức: VD
Một vật chịu hợp lực $0\,\text{N}$. Theo định luật $I$ Newton, gia tốc của vật bằng bao nhiêu (kết quả theo đơn vị m/s^2).
\choice
{\True 0}
{1}
{0}
{0}
\loigiai{
Áp dụng hệ thức phù hợp của bài học, thay số được $0\,\text{m}$ /s^2.
}
\end{ex}

% ===== Câu 145 =====
% ID: AIQ_L10C3_FULL_0307
% Mức: VD
\begin{ex}
% ID: AIQ_L10C3_FULL_0307
% Mức: VD
Một vật chịu hợp lực $0\,\text{N}$. Theo định luật $I$ Newton, gia tốc của vật bằng bao nhiêu (kết quả theo đơn vị m/s^2).
\choice
{1}
{\True 0}
{0}
{0}
\loigiai{
Áp dụng hệ thức phù hợp của bài học, thay số được $0\,\text{m}$ /s^2.
}
\end{ex}

% ===== Câu 146 =====
% ID: AIQ_L10C3_FULL_0308
% Mức: TH
\begin{ex}
% ID: AIQ_L10C3_FULL_0308
% Mức: TH
Một vật chịu hợp lực $0\,\text{N}$. Theo định luật $I$ Newton, gia tốc của vật bằng bao nhiêu (kết quả theo đơn vị m/s^2).
\choiceTF
{\True Giá trị cần tìm là $0\,\text{m}$ /s^2.}
{Giá trị cần tìm là $1\,\text{m}$ /s^2.}
{\True Phải xác định đầy đủ các lực hoặc đại lượng liên quan trước khi tính.}
{Có thể bỏ qua điều kiện áp dụng của công thức.}
\loigiai{
\begin{itemchoice}
\item (Đúng) Giá trị cần tìm là $0\,\text{m}$ /s^2. (Vì): kết quả $0\,\text{m}$ /s^2. b) Sai. c) Đúng. d) Sai vì phải kiểm tra điều kiện áp dụng
\item (Sai) Giá trị cần tìm là $1\,\text{m}$ /s^2.
\item (Đúng) Phải xác định đầy đủ các lực hoặc đại lượng liên quan trước khi tính.
\item (Sai) Có thể bỏ qua điều kiện áp dụng của công thức.
\end{itemchoice}
}
\end{ex}

% ===== Câu 147 =====
% ID: AIQ_L10C3_FULL_0309
% Mức: TH
\begin{ex}
% ID: AIQ_L10C3_FULL_0309
% Mức: TH
Một vật chịu hợp lực $0\,\text{N}$. Theo định luật $I$ Newton, gia tốc của vật bằng bao nhiêu (kết quả theo đơn vị m/s^2).
\choiceTF
{\True Giá trị cần tìm là $0\,\text{m}$ /s^2.}
{Giá trị cần tìm là $1\,\text{m}$ /s^2.}
{\True Phải xác định đầy đủ các lực hoặc đại lượng liên quan trước khi tính.}
{Có thể bỏ qua điều kiện áp dụng của công thức.}
\loigiai{
\begin{itemchoice}
\item (Đúng) Giá trị cần tìm là $0\,\text{m}$ /s^2. (Vì): kết quả $0\,\text{m}$ /s^2. b) Sai. c) Đúng. d) Sai vì phải kiểm tra điều kiện áp dụng
\item (Sai) Giá trị cần tìm là $1\,\text{m}$ /s^2.
\item (Đúng) Phải xác định đầy đủ các lực hoặc đại lượng liên quan trước khi tính.
\item (Sai) Có thể bỏ qua điều kiện áp dụng của công thức.
\end{itemchoice}
}
\end{ex}

% ===== Câu 148 =====
% ID: AIQ_L10C3_FULL_0310
% Mức: VD
\begin{ex}
% ID: AIQ_L10C3_FULL_0310
% Mức: VD
Một vật chịu hợp lực $0\,\text{N}$. Theo định luật $I$ Newton, gia tốc của vật bằng bao nhiêu (kết quả theo đơn vị m/s^2).
\choiceTF
{\True Giá trị cần tìm là $0\,\text{m}$ /s^2.}
{Giá trị cần tìm là $1\,\text{m}$ /s^2.}
{\True Phải xác định đầy đủ các lực hoặc đại lượng liên quan trước khi tính.}
{Có thể bỏ qua điều kiện áp dụng của công thức.}
\loigiai{
\begin{itemchoice}
\item (Đúng) Giá trị cần tìm là $0\,\text{m}$ /s^2. (Vì): kết quả $0\,\text{m}$ /s^2. b) Sai. c) Đúng. d) Sai vì phải kiểm tra điều kiện áp dụng
\item (Sai) Giá trị cần tìm là $1\,\text{m}$ /s^2.
\item (Đúng) Phải xác định đầy đủ các lực hoặc đại lượng liên quan trước khi tính.
\item (Sai) Có thể bỏ qua điều kiện áp dụng của công thức.
\end{itemchoice}
}
\end{ex}

% ===== Câu 149 =====
% ID: AIQ_L10C3_FULL_0311
% Mức: VD
\begin{ex}
% ID: AIQ_L10C3_FULL_0311
% Mức: VD
Một vật chịu hợp lực $0\,\text{N}$. Theo định luật $I$ Newton, gia tốc của vật bằng bao nhiêu (kết quả theo đơn vị m/s^2).
\choiceTF
{\True Giá trị cần tìm là $0\,\text{m}$ /s^2.}
{Giá trị cần tìm là $1\,\text{m}$ /s^2.}
{\True Phải xác định đầy đủ các lực hoặc đại lượng liên quan trước khi tính.}
{Có thể bỏ qua điều kiện áp dụng của công thức.}
\loigiai{
\begin{itemchoice}
\item (Đúng) Giá trị cần tìm là $0\,\text{m}$ /s^2. (Vì): kết quả $0\,\text{m}$ /s^2. b) Sai. c) Đúng. d) Sai vì phải kiểm tra điều kiện áp dụng
\item (Sai) Giá trị cần tìm là $1\,\text{m}$ /s^2.
\item (Đúng) Phải xác định đầy đủ các lực hoặc đại lượng liên quan trước khi tính.
\item (Sai) Có thể bỏ qua điều kiện áp dụng của công thức.
\end{itemchoice}
}
\end{ex}

% ===== Câu 150 =====
% ID: AIQ_L10C3_FULL_0312
% Mức: VD
\begin{ex}
% ID: AIQ_L10C3_FULL_0312
% Mức: VD
Một vật chịu hợp lực $0\,\text{N}$. Theo định luật $I$ Newton, gia tốc của vật bằng bao nhiêu (kết quả theo đơn vị m/s^2).
\choiceTF
{\True Giá trị cần tìm là $0\,\text{m}$ /s^2.}
{Giá trị cần tìm là $1\,\text{m}$ /s^2.}
{\True Phải xác định đầy đủ các lực hoặc đại lượng liên quan trước khi tính.}
{Có thể bỏ qua điều kiện áp dụng của công thức.}
\loigiai{
\begin{itemchoice}
\item (Đúng) Giá trị cần tìm là $0\,\text{m}$ /s^2. (Vì): kết quả $0\,\text{m}$ /s^2. b) Sai. c) Đúng. d) Sai vì phải kiểm tra điều kiện áp dụng
\item (Sai) Giá trị cần tìm là $1\,\text{m}$ /s^2.
\item (Đúng) Phải xác định đầy đủ các lực hoặc đại lượng liên quan trước khi tính.
\item (Sai) Có thể bỏ qua điều kiện áp dụng của công thức.
\end{itemchoice}
}
\end{ex}

% ===== Câu 151 =====
% ID: AIQ_L10C3_FULL_0313
% Mức: TH
\begin{ex}
% ID: AIQ_L10C3_FULL_0313
% Mức: TH
Một vật chịu hợp lực $0\,\text{N}$. Theo định luật $I$ Newton, gia tốc của vật bằng bao nhiêu (kết quả theo đơn vị m/s^2).
\shortans{0}
\loigiai{
Thay số vào hệ thức của bài học thu được $0\,\text{m}$ /s^2.
}
\end{ex}

% ===== Câu 152 =====
% ID: AIQ_L10C3_FULL_0314
% Mức: TH
\begin{ex}
% ID: AIQ_L10C3_FULL_0314
% Mức: TH
Một vật chịu hợp lực $0\,\text{N}$. Theo định luật $I$ Newton, gia tốc của vật bằng bao nhiêu (kết quả theo đơn vị m/s^2).
\shortans{0}
\loigiai{
Thay số vào hệ thức của bài học thu được $0\,\text{m}$ /s^2.
}
\end{ex}

% ===== Câu 153 =====
% ID: AIQ_L10C3_FULL_0315
% Mức: TH
\begin{ex}
% ID: AIQ_L10C3_FULL_0315
% Mức: TH
Một vật chịu hợp lực $0\,\text{N}$. Theo định luật $I$ Newton, gia tốc của vật bằng bao nhiêu (kết quả theo đơn vị m/s^2).
\shortans{0}
\loigiai{
Thay số vào hệ thức của bài học thu được $0\,\text{m}$ /s^2.
}
\end{ex}

% ===== Câu 154 =====
% ID: AIQ_L10C3_FULL_0316
% Mức: VD
\begin{ex}
% ID: AIQ_L10C3_FULL_0316
% Mức: VD
Một vật chịu hợp lực $0\,\text{N}$. Theo định luật $I$ Newton, gia tốc của vật bằng bao nhiêu (kết quả theo đơn vị m/s^2).
\shortans{0}
\loigiai{
Thay số vào hệ thức của bài học thu được $0\,\text{m}$ /s^2.
}
\end{ex}

% ===== Câu 155 =====
% ID: AIQ_L10C3_FULL_0317
% Mức: VD
\begin{ex}
% ID: AIQ_L10C3_FULL_0317
% Mức: VD
Một vật chịu hợp lực $0\,\text{N}$. Theo định luật $I$ Newton, gia tốc của vật bằng bao nhiêu (kết quả theo đơn vị m/s^2).
\shortans{0}
\loigiai{
Thay số vào hệ thức của bài học thu được $0\,\text{m}$ /s^2.
}
\end{ex}

% ===== Câu 156 =====
% ID: AIQ_L10C3_FULL_0318
% Mức: VD
\begin{ex}
% ID: AIQ_L10C3_FULL_0318
% Mức: VD
Một vật chịu hợp lực $0\,\text{N}$. Theo định luật $I$ Newton, gia tốc của vật bằng bao nhiêu (kết quả theo đơn vị m/s^2).
\shortans{0}
\loigiai{
Thay số vào hệ thức của bài học thu được $0\,\text{m}$ /s^2.
}
\end{ex}

% ===== Câu 157 =====
% ID: AIQ_L10C3_FULL_0319
% Mức: VD
\begin{bt}
% ID: AIQ_L10C3_FULL_0319
% Mức: VD
Một vật chịu hợp lực $0\,\text{N}$. Theo định luật $I$ Newton, gia tốc của vật bằng bao nhiêu (kết quả theo đơn vị m/s^2).
\loigiai{
Xác định các đại lượng đã cho, chọn hệ thức phù hợp và thay số. Kết quả: $0\,\text{m}$ /s^2.
}
\end{bt}

% ===== Câu 158 =====
% ID: AIQ_L10C3_FULL_0320
% Mức: VD
\begin{bt}
% ID: AIQ_L10C3_FULL_0320
% Mức: VD
Một vật chịu hợp lực $0\,\text{N}$. Theo định luật $I$ Newton, gia tốc của vật bằng bao nhiêu (kết quả theo đơn vị m/s^2).
\loigiai{
Xác định các đại lượng đã cho, chọn hệ thức phù hợp và thay số. Kết quả: $0\,\text{m}$ /s^2.
}
\end{bt}

% ===== Câu 159 =====
% ID: AIQ_L10C3_FULL_0321
% Mức: VD
\begin{bt}
% ID: AIQ_L10C3_FULL_0321
% Mức: VD
Một vật chịu hợp lực $0\,\text{N}$. Theo định luật $I$ Newton, gia tốc của vật bằng bao nhiêu (kết quả theo đơn vị m/s^2).
\loigiai{
Xác định các đại lượng đã cho, chọn hệ thức phù hợp và thay số. Kết quả: $0\,\text{m}$ /s^2.
}
\end{bt}

% ===== Câu 160 =====
% ID: AIQ_L10C3_FULL_0322
% Mức: VDC
\begin{bt}
% ID: AIQ_L10C3_FULL_0322
% Mức: VDC
Một vật chịu hợp lực $0\,\text{N}$. Theo định luật $I$ Newton, gia tốc của vật bằng bao nhiêu (kết quả theo đơn vị m/s^2).
\loigiai{
Xác định các đại lượng đã cho, chọn hệ thức phù hợp và thay số. Kết quả: $0\,\text{m}$ /s^2.
}
\end{bt}

% ===== Câu 161 =====
% ID: AIQ_L10C3_FULL_0323
% Mức: VDC
\begin{bt}
% ID: AIQ_L10C3_FULL_0323
% Mức: VDC
Một vật chịu hợp lực $0\,\text{N}$. Theo định luật $I$ Newton, gia tốc của vật bằng bao nhiêu (kết quả theo đơn vị m/s^2).
\loigiai{
Xác định các đại lượng đã cho, chọn hệ thức phù hợp và thay số. Kết quả: $0\,\text{m}$ /s^2.
}
\end{bt}

% ===== Câu 162 =====
% ID: AIQ_L10C3_FULL_0324
% Mức: VDC
\begin{bt}
% ID: AIQ_L10C3_FULL_0324
% Mức: VDC
Một vật chịu hợp lực $0\,\text{N}$. Theo định luật $I$ Newton, gia tốc của vật bằng bao nhiêu (kết quả theo đơn vị m/s^2).
\loigiai{
Xác định các đại lượng đã cho, chọn hệ thức phù hợp và thay số. Kết quả: $0\,\text{m}$ /s^2.
}
\end{bt}
